\chapter{经典芯片、启动与平台演进}

前面已经把最小 CPU 和整机事务分开讲清楚。本章再回到真实芯片：6502、Z80、8051、8086、80386、80486、Pentium 和现代 SoC 不只是历史名称，而是硬件边界逐步扩展的样本。读这些芯片时，不应只背年份、位宽或指令集，而要问：状态保存在哪里，地址怎样形成，外部事务怎样发起，异常和等待怎样改变下一步，哪些功能被逐步搬进芯片内部。

\begin{longtable}{L{0.19\linewidth}L{0.29\linewidth}L{0.42\linewidth}}
\toprule
芯片 & 本章观察重点 & 对整机理解的贡献 \\
\midrule
6502/Z80 & 8-bit CPU、外部地址/数据总线、控制周期 & 看见 CPU、存储器和 I/O 如何靠板级总线协作 \\
8051 & CPU、片内存储、I/O、定时器和中断集成 & 看见单片机把一部分整机控制收进芯片 \\
8086/8088 & 16-bit x86、20-bit 地址、复用总线和预取队列 & 看见早期微处理器如何在引脚和地址空间之间取舍 \\
80386 & 32-bit、保护模式、分页和异常机制 & 看见 CPU 如何承担系统级地址和权限约定 \\
486/Pentium & 片上 cache、流水线、浮点和预测方向 & 看见低时延来自片上集成和执行结构重组 \\
现代 SoC & CPU、cache、内存控制器、I/O、安全启动同片集成 & 看见整机边界继续向片内移动 \\
\bottomrule
\end{longtable}

\begin{classicnote}
本章采用经典系统教材常用的标注方式：不把 CPU、总线、cache、I/O 当作孤立名词，而把它们写成可追踪的事务。阅读任何一颗芯片时，都应同时回答五个问题：\hwkey{状态}保存在哪里，\hwkey{组合路径}经过哪些部件，\hwkey{控制信号}在什么时候有效，\hwkey{外部事务}由谁驱动和谁采样，\hwkey{异常或等待}如何改变下一步。
\end{classicnote}

\begin{pdfdiagram}
  \node[hw lane, minimum width=1.6cm] (a) at (0,0) {\shortstack{8-bit\\外部总线}};
  \node[hw lane, minimum width=1.65cm] (b) at (2.2,0) {\shortstack{8086\\复用总线}};
  \node[hw lane, minimum width=1.75cm] (c) at (4.55,0) {\shortstack{80386\\保护/分页}};
  \node[hw lane, minimum width=1.75cm] (d) at (6.95,0) {\shortstack{486/Pentium\\cache/流水线}};
  \node[hw lane, minimum width=1.7cm] (e) at (9.25,0) {\shortstack{SoC\\片上平台}};
  \draw[hw bus] (a) -- (b);
  \draw[hw bus] (b) -- (c);
  \draw[hw bus] (c) -- (d);
  \draw[hw bus] (d) -- (e);
  \node[hw label, text width=9.4cm] at (4.65,-1.0) {演进主线不是“位宽越来越大”这么简单，而是地址形成、等待隐藏、异常诊断和外设事务逐步进入芯片内部。};
\end{pdfdiagram}
\diagramnote{经典芯片观察线索。每一代都把前一章的事务边界移动、扩宽或缓存。}

\section{一、早期微处理器：从 6502、Z80 到 8086}

在 8086 之前，8-bit 微处理器已经把“CPU 加外部总线”这台戏完整演过一遍。6502、Z80 和 8051 各自回答了一个整机问题：寄存器太少怎么办，上下文切换和 DRAM 刷新如何省掉总线事务，整机部件能不能收进同一颗芯片。

\topic{6502 用极少的寄存器和零页换取简单的数据通路}6502 的编程模型只有累加器 A、变址寄存器 X 和 Y、栈指针 SP、程序计数器 PC 和标志寄存器 P，没有一组可自由选用的通用寄存器；整颗芯片约 3500--4500 个晶体管，数据通路围绕累加器组织，没有乘除法指令。寄存器这样少，工作数据就必须经常放在存储器里，6502 的对策是把地址空间最低的第 0 页（\code{0x0000}--\code{0x00FF}）当作“伪寄存器堆”使用：零页寻址的指令只有 2 个字节，3 个时钟拍完成一次读写，比访问任意 16-bit 地址更短更快。栈则被硬件固定在第 1 页（\code{0x0100}--\code{0x01FF}），SP 因此只需保存 8-bit 页内偏移。数组处理依靠 \code{(zp),Y} 间接变址：从零页的一对字节取出 16-bit 基地址，再加上 Y 作为元素偏移，一次读用 5 拍，相加后若跨页再加 1 拍。

\begin{longtable}{L{0.2\linewidth}L{0.12\linewidth}L{0.58\linewidth}}
\toprule
指令形式 & 拍数 & 对应的总线动作 \\
\midrule
LDA 立即数 & 2 & 取操作码，再把操作数字节直接读入 A \\
LDA 零页 & 3 & 取操作码，取零页地址，从第 0 页读数据 \\
LDA 绝对地址 & 4 & 取操作码，取 16-bit 地址的两个字节，再读数据 \\
LDA \code{(zp),Y} & 5+1 & 取操作码，从零页读基地址两字节，加 Y 后读数据，跨页加 1 拍 \\
STA \code{(zp),Y} & 6 & 取操作码，从零页读基地址两字节，加 Y 形成地址后写入，写比读多一拍 \\
JMP 绝对地址 & 3 & 取操作码，取目标地址两个字节并装入 PC \\
\bottomrule
\end{longtable}

每条指令的拍数固定而且很小，总线在每个时钟拍上完成一个确定动作，因此 6502 程序可以手工逐拍 trace（执行轨迹）：哪一拍取指、哪一拍读零页、哪一拍形成地址，全能数清楚。6502 的便宜正来自这种简单：没有乘除法，没有多通用寄存器，数据通路只围绕累加器，芯片面积和售价才可能压到最低。

\topic{Z80 用双寄存器组和取指刷新换取系统效率}Z80 在软件上与 8080 兼容，并在此之上做了几处扩展。最显眼的是寄存器组：主寄存器组 AF、BC、DE、HL 之外还有整套备用组 AF'、BC'、DE'、HL'，\code{EX} 和 \code{EXX} 各用一条指令完成整组交换。中断响应时不必把寄存器逐个压栈再逐个恢复，现场切换不发生任何存储器访问。第二个扩展与 DRAM 刷新有关：R 寄存器保存刷新行地址，在每个 M1 取指周期的后半段——此时指令字节已经读入、地址总线正好空闲——自动把 R 输出到地址总线上完成一次刷新，随后自动加一；刷新被“借”进取指周期，外部不再需要单独的刷新周期。此外，I 寄存器配合中断模式 2 提供向量页地址，IX、IY 变址寄存器扩展了寻址方式。

\begin{longtable}{L{0.22\linewidth}L{0.4\linewidth}L{0.28\linewidth}}
\toprule
Z80 机制 & 硬件做法 & 换取的收益 \\
\midrule
主/备寄存器组 & AF/BC/DE/HL 与 AF'/BC'/DE'/HL'，\code{EX}/\code{EXX} 整组交换 & 中断现场切换不访问存储器 \\
R 寄存器 & M1 取指后半段输出刷新地址，随后自动加一 & DRAM 刷新不占额外总线周期 \\
I 寄存器与中断模式 2 & I 提供向量页地址，设备提供页内偏移 & 中断入口可以成表组织 \\
IX/IY 变址寄存器 & 变址基址加偏移形成有效地址 & 数据结构和栈帧访问更直接 \\
\bottomrule
\end{longtable}

Z80 的教学点有两条。第一，寄存器越多，上下文切换越少访问内存：一组备用寄存器换来的是中断响应里整段存储器往返的消失。第二，刷新是“借用已有事务”的设计：DRAM 刷新本来要占用专门的总线时间，Z80 把它塞进每个取指周期里必然出现的空闲半段，系统里就再也看不见它。

\topic{8051 把一部分整机收进同一颗芯片}8051 是单片机：CPU、4--8 KiB ROM/EPROM、128--256 B 片内 RAM、4 个 8-bit 端口、2--3 个定时器、UART 和中断控制器全部在同一颗芯片内。片内 RAM 分区使用：最低 32 字节是 4 组工作寄存器 R0--R7，由程序状态字 PSW 的 RS0、RS1 两位选组，一条改选组的指令就完成一次快速上下文切换；\code{0x20}--\code{0x2F} 的 16 字节是位寻址区，其中 128 个 bit 可以单独寻址、置位和清零，标志位不必占用整个字节；其余部分是通用区。片内外设寄存器称为特殊功能寄存器（SFR），占用 \code{0x80}--\code{0xFF} 的直接地址。四个端口采用准双向结构：引脚锁存器先写入 1，该引脚才能作为输入被读取。

8051 的教学点直接连着第五章：把 ROM、RAM、地址译码、MMIO、定时器、串口和中断控制器收进一颗芯片之后，“总线”变成了片内连线，板级事务缩成片内信号，但语义没有改变——写 SFR 仍然是写设备寄存器，读端口仍然是 MMIO 读，定时器溢出仍然是一个中断事件。整机边界在这里第一次明显移动，而读者已经学过的所有约定原样保留。

8086 是理解早期 x86 体系的重要起点。它是 16-bit 微处理器，拥有 16-bit 通用寄存器和 16-bit 数据通路特征，同时通过 20-bit 地址能力访问最多 1 MiB 物理地址空间。它的历史意义不只在指令集，还在于展示了 CPU 芯片如何通过外部总线、地址锁存、控制信号和存储器/外设共同构成一台机器。

8086 的寄存器组包括 AX、BX、CX、DX、SP、BP、SI、DI，以及 CS、DS、SS、ES 等段寄存器。段寄存器和偏移地址共同形成物理地址：

\begin{lstlisting}
physical_address = segment * 16 + offset
\end{lstlisting}

这个规则反映了一个时代的硬件取舍：内部仍以 16-bit 编程模型为主，但外部需要访问超过 64 KiB 的地址空间，于是通过段基址左移和偏移相加得到 20-bit 地址。分段不是高级抽象，而是地址形成硬件的一部分。

8086 的外部总线也值得注意。地址线和数据线存在复用，某些引脚在总线周期的不同阶段承担不同含义。硬件需要地址锁存器在地址有效阶段保存地址，随后这些引脚可以用于数据传输。一次总线访问不是“读内存”三个字，而是分阶段发生的事务：输出地址、锁存地址、给出读写控制、等待存储器或 I/O 响应、传输数据、结束总线周期。

\begin{longtable}{L{0.24\linewidth}L{0.58\linewidth}}
\toprule
8086 观察点 & 硬件含义 \\
\midrule
16-bit 寄存器 & 数据运算和编程模型以 16-bit 为核心 \\
20-bit 地址形成 & 通过段寄存器和偏移访问 1 MiB 物理空间 \\
地址/数据复用总线 & 芯片引脚有限，外部需要锁存和分阶段传输 \\
存储器与 I/O 周期 & 访问目标由控制信号和地址空间规则区分 \\
BIU/EU 分工 & 取指队列和执行部件可部分重叠工作 \\
\bottomrule
\end{longtable}

8086 的典型板级连接还需要外部支持器件。地址/数据复用引脚在 T1 阶段输出地址，在后续阶段承载数据，因此常用地址锁存器保存低位地址；时钟、复位和 ready 同步可由时钟发生器组织；最大模式下，总线控制信号还可由总线控制器从状态信号译出。读数据手册时，不应只看“8086 是 16-bit CPU”，还要看它把哪些系统责任交给了板级芯片。

\begin{longtable}{L{0.2\linewidth}L{0.34\linewidth}L{0.36\linewidth}}
\toprule
信号或器件 & 硬件作用 & 关键问题 \\
\midrule
ALE & 地址锁存允许，提示外部锁存低位地址 & 锁存器是否在地址有效窗口保存地址 \\
AD0--AD15 & 复用地址/数据线 & 哪个阶段由 CPU 驱动，哪个阶段由存储器或 I/O 驱动 \\
A16--A19/S3--S6 & 高位地址与状态复用 & 地址是否在访问开始阶段被正确保存 \\
\code{RD}/\code{WR} & 读写控制 & 目标芯片输出使能和写使能是否互斥 \\
M/IO & 区分存储器周期和 I/O 周期 & 地址译码是否把事务送到正确目标 \\
READY & 外部目标请求插入等待状态 & 慢存储器是否能可靠延长总线周期 \\
INTR/NMI & 可屏蔽和不可屏蔽中断入口 & 请求是否同步，响应和清除顺序是否明确 \\
8284/8288 & 时钟/复位/ready 或总线控制辅助 & 外部控制信号是否满足时序和极性要求 \\
\bottomrule
\end{longtable}

\topic{读总线周期要按 T 状态逐段读}上一张信号表说明了每个引脚负责什么，但还不够：同一个总线周期里，同一组引脚在不同时间段做不同的事。8086 把一次总线访问切成 T1、T2、T3、T4 四个状态，必要时在 T3 与 T4 之间插入等待状态 Tw。读手册的正确方式是逐状态确认：这一拍谁在驱动地址/数据引脚，哪个控制信号有效，外部芯片应在哪个边沿锁存或采样。

\begin{longtable}{L{0.1\linewidth}L{0.28\linewidth}L{0.24\linewidth}L{0.28\linewidth}}
\toprule
T 状态 & 地址/数据引脚 & 控制信号 & 外部动作 \\
\midrule
T1 & CPU 在 AD0--AD15 和 A16--A19/S3--S6 上输出地址 & ALE 拉高 & 地址锁存器在 ALE 下降沿保存地址 \\
T2 & 地址从复用脚撤出，引脚准备转入数据方向 & \code{RD}（读）或 \code{WR}（写）有效 & 数据收发器按读/写确定传输方向 \\
T3 & 读周期中由目标芯片驱动数据线 & 控制保持，READY 为低则插入 Tw & 存储器或 I/O 把数据放到总线上，等待状态期间地址和控制保持不变 \\
T4 & CPU 采样读入的数据，或结束写周期 & 控制信号撤销 & 目标释放数据线，总线周期结束 \\
\bottomrule
\end{longtable}

写周期与读周期骨架相同，差别在数据方向：CPU 从 T2 起亲自驱动数据线，并保持到写窗口结束，目标芯片在写控制有效的窗口内采样数据。“读总线周期”读到最后，就是把每个 T 状态里谁驱动什么、谁采样什么写清楚；有了这张表，上一张信号表里的每个信号才有了时间轴上的位置。

\topic{8086 的存储体选择是理解 16-bit 总线的关键}8086 有 16-bit 数据总线，但外部存储器常按 8-bit 芯片组织成偶地址低字节存储体和奇地址高字节存储体。地址最低位 A0 与高字节使能 $\overline{\mathrm{BHE}}$ 共同决定哪些字节参与本次传输。这样，16-bit 字访问在偶地址上可以一次完成；若字从奇地址开始，通常需要拆成两个总线周期。

\begin{longtable}{L{0.14\linewidth}L{0.18\linewidth}L{0.28\linewidth}L{0.3\linewidth}}
\toprule
A0 & $\overline{\mathrm{BHE}}$ & 参与字节 & 典型含义 \\
\midrule
0 & 0 & D0--D7 与 D8--D15 & 偶地址 16-bit 字访问 \\
0 & 1 & D0--D7 & 偶地址低字节访问 \\
1 & 0 & D8--D15 & 奇地址字节访问 \\
1 & 1 & 无有效字节 & 正常访问中不应作为有效传输 \\
\bottomrule
\end{longtable}

这张表能解释为什么“16-bit CPU”不等于每次总线访问都简单搬 16 bit。存储芯片、地址译码器和数据收发器必须理解字节车道。若把两个 8-bit 存储体简单并在一起而不处理 A0 和 $\overline{\mathrm{BHE}}$，字节写入会破坏相邻数据，奇地址字访问也无法正确完成。

\topic{最小模式、最大模式和外部控制器}8086 可以在最小模式下直接输出较多总线控制信号，也可以在最大模式下把状态信息交给外部总线控制器生成控制信号。前者适合单处理器小系统，后者便于接入更复杂的总线环境。这个差别说明，CPU 引脚不是固定地“表示某个概念”，它们会随系统模式承担不同职责。

\begin{longtable}{L{0.18\linewidth}L{0.34\linewidth}L{0.38\linewidth}}
\toprule
系统形态 & 典型外部对象 & 学习重点 \\
\midrule
最小模式 & 地址锁存器、数据收发器、ROM/RAM、简单 I/O & CPU 直接给出读写、方向、使能等控制 \\
最大模式 & 8288 总线控制器、总线仲裁、协处理或多主设备 & 状态信号被外部控制器译成总线命令 \\
8088 小系统 & 8-bit 外部数据总线、较低成本板级存储器 & 编程模型相近，外部总线带宽和周期不同 \\
\bottomrule
\end{longtable}

一个 8086 存储器读周期可以按 T 状态走读。T1 中，CPU 把地址放到复用总线并拉出 ALE，外部锁存器保存地址；T2 中，CPU 改变总线方向并给出读控制；T3/Tw 中，存储器或 I/O 设备把数据放到数据线上，若 READY 不满足，CPU 插入等待状态；T4 中，CPU 采样数据并结束周期。这个过程把第一章的三态总线、第二章的时钟边界和第三章的控制 trace（执行轨迹）合到一起：总线上的每一段时间都必须有唯一驱动者和明确采样者。

\begin{lstlisting}
8086 memory read sketch:
  T1: address valid, ALE latches address
  T2: RD active, bus turns around for target data
  T3: target drives data, READY may insert wait states
  T4: CPU samples data, control deasserts
\end{lstlisting}

\begin{pdfdiagram}
  \node[hw lane, minimum width=1.45cm] (t1) at (0,0) {\shortstack{T1\\地址}};
  \node[hw lane, minimum width=1.55cm] (t2) at (2.1,0) {\shortstack{T2\\转向}};
  \node[hw lane, minimum width=1.65cm] (tw) at (4.35,0) {\shortstack{T3/Tw\\等待}};
  \node[hw lane, minimum width=1.45cm] (t4) at (6.65,0) {\shortstack{T4\\采样}};
  \draw[hw bus] (t1) -- node[above=2pt,hw label,fill=white,inner sep=1pt]{ALE 锁存} (t2);
  \draw[hw bus] (t2) -- node[above=2pt,hw label,fill=white,inner sep=1pt]{RD 有效} (tw);
  \draw[hw return] (tw) -- node[above=2pt,hw label,fill=white,inner sep=1pt]{READY} (t4);
  \node[hw label, text width=2.6cm] at (0,-1.0) {CPU 驱动地址};
  \node[hw label, text width=2.8cm] at (2.1,-1.0) {总线方向切换};
  \node[hw label, text width=3.1cm] at (4.35,-1.0) {目标驱动数据};
  \node[hw label, text width=2.8cm] at (6.65,-1.0) {CPU 采样数据};
  \node[hw label, text width=9.3cm] at (3.35,-2.0) {复用总线不是一根线同时表示所有含义，而是在同一个读周期内按 T 状态依次承担地址、方向、数据和完成边界。};
\end{pdfdiagram}
\diagramnote{8086 存储器读周期。地址锁存、总线转向、等待和采样必须分阶段观察。}

\topic{8086 写周期、中断响应和总线让出也要按事务读}读周期只说明目标驱动数据线；写周期则相反，CPU 在合适阶段驱动数据线，外部存储器或 I/O 设备在写控制有效窗口内采样。若数据收发器方向或写使能时序错误，写入可能丢失，也可能造成 CPU 和外设同时驱动总线。8086 的中断响应周期又是另一类事务：CPU 不访问普通内存数据，而是响应中断控制器，取得中断类型或进入规定入口。

\begin{longtable}{L{0.18\linewidth}L{0.36\linewidth}L{0.36\linewidth}}
\toprule
事务 & 总线动作 & 观察重点 \\
\midrule
存储器写 & 地址先被锁存，CPU 驱动数据，\code{WR} 有效，目标采样 & 数据方向、字节存储体和写窗口 \\
I/O 写 & 地址或端口号选择外设寄存器，CPU 驱动命令或数据 & M/IO 或 I/O 周期控制、端口译码 \\
中断响应 & CPU 接受 INTR 后执行响应周期，中断控制器给出类型或向量信息 & 请求保持、优先级、响应和屏蔽状态 \\
DMA 让总线 & DMA 请求总线，CPU 完成当前边界后让出，总线主控权转移 & HOLD/HLDA 类握手、三态释放、仲裁 \\
\bottomrule
\end{longtable}

\begin{lstlisting}
8086 memory write sketch:
  T1: address valid, ALE latches address
  T2: CPU drives write data, WR becomes active
  T3: target samples data, READY may insert wait states
  T4: WR deasserts, bus returns to idle or next cycle
\end{lstlisting}

这组事务可以直接指导早期微机板级排错。若读 ROM 正常、写 RAM 失败，应重点查数据收发器方向、\code{WR} 极性、字节存储体选择和 SRAM 写入时间；若外设中断不进入，应同时查外设请求、8259 屏蔽位、CPU 中断允许状态和中断响应周期；若 DMA 后内存内容异常，应查总线让出时 CPU 是否真正三态、DMA 地址计数是否正确、目标存储器是否满足时序。

8086 内部常被划分为总线接口单元和执行单元。总线接口单元负责取指、形成物理地址、访问外部总线，并维护预取队列；执行单元负责译码和执行指令。这个分工说明，即使没有现代深流水，CPU 内部也已经在尝试把“取指访问外部世界”和“执行内部运算”分开，以减少外部存储访问延迟对执行的影响。

从本书前面建立的硬件模型看，8086 可以视为一个更复杂的同步系统：PC 在 x86 语境中体现为指令指针和代码段组合；地址形成部件把段和偏移送入加法路径；总线接口把地址和控制信号变成外部事务；执行部件译码指令并驱动寄存器、ALU 和标志位。8086 不是脱离前面章节的新对象，而是前述部件在真实芯片中的组合。

8086 还说明“低时延”在早期处理器中已经是系统问题。预取队列试图在执行部件工作时提前从外部总线取回后续指令，减少执行部件等待取指的时间。这个机制不等于现代流水线和分支预测，但思想相通：把慢的外部访问与内部执行尽量重叠。若分支改变控制流，预取队列中已经取得的顺序指令可能失效；若外部总线慢于指令消费速度，执行部件仍会等待。这里可以看见第一章快慢路径表的早期版本。

\begin{chipcase}
围绕 8086/8088 的经典外围芯片同样值得提前认识。8255 可编程并行接口把若干引脚组织成可配置 I/O 口；8253/8254 定时器提供周期性计数和中断来源；8259 中断控制器把多个外设请求汇总给 CPU；8237 DMA 控制器让设备和内存之间进行批量搬运；8250/16450/16550 UART 把串行线上的位流转换为 CPU 可读写的寄存器和 FIFO。它们共同说明，整机不是一颗 CPU，而是一组围绕地址、数据、控制和完成事件协作的芯片。
\end{chipcase}

\begin{longtable}{L{0.18\linewidth}L{0.34\linewidth}L{0.38\linewidth}}
\toprule
经典芯片 & 解决的问题 & 与本章主题的连接 \\
\midrule
8255 PPI & 把并行引脚配置成输入、输出或握手端口 & GPIO 和设备寄存器的早期形态 \\
8253/8254 PIT & 用计数器产生周期、延时和中断节拍 & 定时器寄存器与完成事件 \\
8259 PIC & 汇总、屏蔽和优先级排序中断请求 & 中断控制器和 CPU 入口 \\
8237 DMA & 在设备和内存之间搬运数据 & 总线主设备、仲裁和 cache 可见性 \\
8250/16550 UART & 串行收发、状态位和 FIFO & 设备状态、数据寄存器和中断通知 \\
\bottomrule
\end{longtable}

\topic{8088/8086 最小系统可以按模块逐步实现}若要把 8086 类系统搭成现实电路，不应从“跑 DOS”开始，而应从最小可运行系统开始。8088 常用于教学，因为外部数据总线为 8 bit，存储器接口更容易搭建；8086 则保留 16-bit 外部数据总线，需要同时处理奇偶字节存储体。无论选择哪一种，系统都要包括时钟/复位、地址锁存、数据收发、ROM、RAM、I/O 译码和至少一个可观察外设。

\begin{longtable}{L{0.2\linewidth}L{0.32\linewidth}L{0.38\linewidth}}
\toprule
模块 & 典型器件或做法 & 第一轮检查 \\
\midrule
时钟/复位 & 8284 或等价时钟复位电路 & 复位保持期间写控制无效，释放后第一拍地址可预测 \\
地址锁存 & 74LS373/8282 锁存 AD 线上低位地址 & ALE 有效窗口内保存地址，数据阶段地址不丢失 \\
数据收发 & 74LS245/8286 控制数据方向 & 读周期外设驱动，写周期 CPU 驱动，不发生总线争用 \\
启动 ROM & EPROM/Flash 映射到复位入口 & 复位入口读出跳转或初始化代码 \\
SRAM & 静态 RAM 与字节使能/写使能连接 & 写一个地址后再读回同一数据，邻近字节不变 \\
I/O 译码 & 74LS138、比较器或 PAL/GAL & 8255、8253、8259、UART 等端口互不重叠 \\
可观察外设 & LED 锁存器、按键、串口或逻辑分析仪探针 & CPU 写端口后能看到确定副作用 \\
\bottomrule
\end{longtable}

这样的系统可以分三步 bring-up。第一步只让 CPU 从 ROM 取指，确认复位入口、地址锁存和读周期成立；第二步加入 SRAM，执行一个写后读回的小程序，确认写周期、数据方向和等待状态成立；第三步加入 8255 或简单 LED 锁存器，确认 I/O 周期、端口译码和外设副作用成立。每一步都应记录地址线、ALE、读写控制、数据总线方向和目标片选。若第一步没有通过，后续软件调试没有意义。

\topic{等待状态发生器是慢器件和 CPU 之间的约定}8086/8088 的 READY 信号把固定周期总线变成可延长事务。慢 ROM、慢 SRAM、外设端口或板级译码较深时，目标可以通过 READY 让 CPU 插入等待状态。等待状态不是“让程序慢一点”的软件概念，而是总线周期中的硬件暂停：地址、控制和方向必须保持可解释，目标在满足访问时间后再允许 CPU 采样或结束写周期。

\begin{longtable}{L{0.18\linewidth}L{0.34\linewidth}L{0.38\linewidth}}
\toprule
对象 & READY/等待的作用 & 关键问题 \\
\midrule
慢 ROM & 延长取指或读常量周期 & 等待期间地址和片选是否保持稳定 \\
慢 SRAM & 延长数据读写窗口 & 写周期中数据是否覆盖完整写入窗口 \\
I/O 外设 & 给外设状态机足够时间返回数据 & 外设未 ready 时 CPU 是否错误采样浮空总线 \\
DMA/仲裁 & 让总线所有权变化落在安全边界 & CPU 是否在让总线前完成当前事务 \\
\bottomrule
\end{longtable}

一个简单等待状态发生器可以由地址译码、访问类型和计数器组成：快 SRAM 区域不插等待，慢 ROM 区域插入一个等待，外设区域插入两个等待或等待外设 ACK。正式规格应说明等待的最大长度和超时行为。否则一个坏外设可能永久拉住 CPU，使整机看起来像死机。

\topic{PC/XT 到 PC/AT 是早期整机平台案例}8088、8086、80286 这些处理器进入实际计算机后，通常不会单独出现。以 PC/XT、PC/AT 风格的平台为例，CPU 周围有启动 ROM、DRAM、DMA 控制器、中断控制器、定时器、键盘控制、串口、并口、显示适配器和扩展总线。它们共同定义了“软件看到的机器”。因此，理解这类平台时，不能只问 CPU 支持哪些指令，还要问端口地址怎样分配、中断线怎样汇总、DMA 通道怎样仲裁、启动 ROM 映射在哪里、扩展卡如何把自己的寄存器和 ROM 挂进系统。

\begin{longtable}{L{0.18\linewidth}L{0.34\linewidth}L{0.38\linewidth}}
\toprule
平台对象 & 硬件职责 & 预期行为 \\
\midrule
BIOS ROM & 复位后提供第一段固件，初始化基本设备 & 复位入口命中 ROM，高地址别名或跳转正确 \\
8259A PIC & 汇总外设中断并向 CPU 投递 & 屏蔽位、优先级、EOI 和中断向量正确 \\
8253/8254 PIT & 提供周期节拍和计数事件 & 计数输入、模式字、输出和中断请求正确 \\
8237A DMA & 在设备和内存之间搬运块数据 & 通道、页寄存器、地址计数和终端计数正确 \\
键盘控制器 & 接收键盘扫描码，也常参与复位/门控控制 & 输入缓冲、输出缓冲和命令应答顺序正确 \\
UART/并口 & 提供串行或并行外设连接 & 状态位、FIFO 或握手线能解释数据何时可收发 \\
扩展总线 & 允许插卡响应 I/O、内存窗口或中断线 & 片选互斥、总线方向、中断线和 DMA 请求不冲突 \\
\bottomrule
\end{longtable}

这个案例把本章许多概念放到同一台机器里。BIOS ROM 是启动约定；8259 是中断入口；8254 是定时事件；8237 是 DMA 总线主设备；UART 是状态寄存器与数据寄存器；扩展总线是地址译码、三态和仲裁的板级版本。若一台 PC/AT 风格机器无法启动，应按硬件路径排查：复位后 CPU 是否访问 BIOS 入口，ROM 是否驱动数据总线，时钟和 READY 是否满足，随后才检查键盘、显示、磁盘或操作系统。

\topic{端口地址表是平台规格的一部分}早期 x86 平台大量使用独立 I/O 空间。设备寄存器不一定映射在普通内存地址中，而是通过端口号和 I/O 周期访问。平台必须写清楚哪些端口归哪个设备，哪些中断线归哪个设备，哪些 DMA 通道归哪个设备。否则，即使 CPU 指令正确，两个设备也可能响应同一个端口，或一个中断请求被错误解释成另一个设备事件。

\begin{longtable}{L{0.2\linewidth}L{0.32\linewidth}L{0.38\linewidth}}
\toprule
规格项 & 说明内容 & 典型错误 \\
\midrule
I/O 端口 & 控制、状态、数据寄存器的端口号和读写语义 & 端口译码过宽，多个设备同时响应 \\
中断线 & 哪个设备接到哪条 IRQ，是否可屏蔽和共享 & 设备已请求但 PIC 屏蔽，或 EOI 顺序错误 \\
DMA 通道 & 通道号、方向、地址和长度寄存器 & 8-bit/16-bit 通道混用或页寄存器错误 \\
内存窗口 & 显存、Option ROM、缓冲区或 MMIO 区域 & RAM、ROM 和设备窗口重叠 \\
启动顺序 & 固件先初始化哪些控制器，再枚举哪些设备 & 设备尚未复位稳定就被访问 \\
\bottomrule
\end{longtable}

把端口地址表写成规格，有助于理解“平台兼容性”。处理器可能更换为更快的 80286、80386 或 80486，但许多 I/O 端口、中断和启动约定仍然保留，使旧软件能够找到键盘、定时器、串口和磁盘控制器。硬件的发展不是每一代从零开始，而是在保留平台约定的同时，把部分功能移入更高集成度的芯片组。

\section{二、80386：32-bit、保护模式与地址转换}

80386 把 x86 推入 32-bit 时代。它扩大了通用寄存器宽度，引入 EAX、EBX、ECX、EDX、ESP、EBP、ESI、EDI 等 32-bit 形式；地址和数据处理能力显著增强，并提供保护模式、分页支持和更完整的内存保护机制。相较 8086，80386 展示的不只是“位宽变大”，而是 CPU 开始承担更复杂的系统级硬件职责。

80386 的一个关键变化是 32-bit 线性地址。程序形成的有效地址先经过分段机制得到线性地址；在启用分页时，线性地址再经过页表转换为物理地址。可以用简化链路表示：

\begin{lstlisting}
effective address -> segmentation -> linear address
linear address    -> paging       -> physical address
\end{lstlisting}

这条链路说明，地址不再只是“寄存器加偏移后送到总线”。CPU 内部必须有地址形成、权限检查、页表访问、异常触发等硬件机制。若页表项不存在、权限不允许、访问越界，CPU 不是简单得到一个错误数据，而是产生异常事件，控制流转入规定入口。

保护模式把“谁能访问什么”变成硬件可检查的约定。段描述符、特权级、页表权限等机制，使得不同程序或系统组件不能随意访问全部内存或执行全部指令。这些机制后来成为现代操作系统隔离、虚拟内存和系统调用的硬件基础。对本书来说，重点不是展开操作系统，而是看见硬件层面的扩展：CPU 不只执行 ALU 运算，还参与地址翻译、权限判断和异常入口选择。

\begin{longtable}{L{0.24\linewidth}L{0.58\linewidth}}
\toprule
80386 观察点 & 硬件含义 \\
\midrule
32-bit 通用寄存器 & 数据通路和编程模型扩展到 32-bit \\
32-bit 线性地址 & 地址空间显著扩大，地址形成更复杂 \\
保护模式 & CPU 硬件检查权限、界限和特权级 \\
分页机制 & 线性地址可经页表映射到物理地址 \\
异常机制 & 地址或权限错误进入硬件规定的控制入口 \\
\bottomrule
\end{longtable}

80386 的外部接口同样要按事务理解。32-bit 数据路径并不意味着每次访问都天然是完整 32-bit 字。处理器需要用地址低位和字节使能说明本次访问哪些字节有效；外部总线、存储器阵列和 cache 控制器必须据此选择正确字节、处理未对齐访问或拆分周期。若只把 80386 理解成“寄存器变成 32 位”，就会漏掉字节使能、总线周期、异常和地址转换这些整机边界。

\begin{longtable}{L{0.22\linewidth}L{0.34\linewidth}L{0.34\linewidth}}
\toprule
80386 对象 & 硬件含义 & 常见调试点 \\
\midrule
地址线与数据线 & 指出事务目标并传输 32-bit 数据片段 & 地址对齐、字节选择和目标译码是否一致 \\
字节使能 & 表明本周期哪些字节有效 & 字节/半字/字访问不会破坏相邻字节 \\
总线状态 & 区分读、写、取指、确认等周期语义 & 外部控制器是否按状态产生正确控制线 \\
分页部件 & 查页目录和页表，形成物理地址 & 缺页、权限和存在位错误是否触发异常 \\
TLB 思想 & 缓存近期地址转换结果 & 命中降低转换时延，失效和权限改变要刷新 \\
\bottomrule
\end{longtable}

分页可以用一次内存加载来观察。指令先形成有效地址，分段得到线性地址；分页开启时，线性地址被拆成页目录索引、页表索引和页内偏移。若转换结果已在 TLB 中，地址转换可以很快完成；若未命中，就要访问内存中的页表结构。页表项还带有存在位、读写权限、用户/内核权限、访问位和脏位等硬件标志。任何一步失败，这次访问不应返回任意数据，而应进入异常处理路径。

\begin{lstlisting}
80386-style address walk:
  effective address
    -> segmentation checks
    -> linear address
    -> TLB lookup or page-table walk
    -> permission checks
    -> physical address or exception
\end{lstlisting}

\begin{pdfdiagram}
  \node[hw state, minimum width=1.8cm] (ea) at (0,0.6) {有效地址};
  \node[hw compute, minimum width=1.95cm] (seg) at (2.35,0.6) {分段检查};
  \node[hw state, minimum width=1.95cm] (lin) at (4.8,0.6) {线性地址};
  \node[hw compute, minimum width=1.95cm] (tlb) at (7.25,0.6) {TLB/页表};
  \node[hw state, minimum width=1.95cm] (pa) at (9.7,0.6) {物理地址};
  \node[hw control, minimum width=2.2cm] (ex) at (7.25,-1.15) {异常入口};
  \draw[hw bus] (ea) -- node[above=2pt,hw label,fill=white,inner sep=1pt]{选择子/偏移} (seg);
  \draw[hw bus] (seg) -- node[above=2pt,hw label,fill=white,inner sep=1pt]{界限/权限通过} (lin);
  \draw[hw bus] (lin) -- node[above=2pt,hw label,fill=white,inner sep=1pt]{PDE/PTE} (tlb);
  \draw[hw bus] (tlb) -- node[above=2pt,hw label,fill=white,inner sep=1pt]{允许访问} (pa);
  \draw[hw return] (seg.south) |- node[pos=0.45,left,hw label,fill=white,inner sep=1pt]{段错误} (ex.west);
  \draw[hw return] (tlb.south) -- node[right,hw label,fill=white,inner sep=1pt]{页错误} (ex.north);
  \node[hw label, text width=9.5cm] at (4.85,-2.15) {80386 的一次内存加载在真正访问 cache 或外部总线前，已经可能因为段、页或权限检查失败而进入异常路径。};
\end{pdfdiagram}
\diagramnote{80386 地址转换链路。地址转换和权限检查是访存事务的一部分，不是软件注释。}

\topic{用一个线性地址拆开分页硬件}以线性地址 \code{0x00403ABC} 为例，80386 的 4 KiB 页式分页可把它拆成页目录索引、页表索引和页内偏移。高 10 bit 选择页目录项，中间 10 bit 选择页表项，低 12 bit 是页内偏移。该地址的页目录索引为 1，页表索引为 3，页内偏移为 \code{0xABC}。若 CR3 指向的页目录中第 1 项不存在，访问在真正读数据前就会变成异常；若页表项存在但权限不允许，也不能继续把物理地址送到 cache 或外部总线。

\begin{longtable}{L{0.2\linewidth}L{0.24\linewidth}L{0.46\linewidth}}
\toprule
字段 & 示例值 & 硬件用途 \\
\midrule
页目录索引 & 1 & 选择页目录中的一个 PDE \\
页表索引 & 3 & 选择对应页表中的一个 PTE \\
页内偏移 & \code{0xABC} & 与物理页框基址相加形成最终地址 \\
P/RW/US 等位 & 来自 PDE/PTE & 判断存在性、读写权限和用户/特权访问 \\
\bottomrule
\end{longtable}

\topic{控制寄存器和描述符表是系统约定的入口}80386 的保护模式不能只理解为“地址变大”。处理器必须知道段描述符在哪里、异常入口在哪里、页目录在哪里、当前是否启用保护和分页。这些信息由控制寄存器、GDTR/IDTR 等状态保存。它们不是普通程序变量，而是 CPU 执行内存访问、异常和特权检查时自动使用的硬件约定。

\begin{longtable}{L{0.18\linewidth}L{0.32\linewidth}L{0.4\linewidth}}
\toprule
对象 & 保存内容 & 硬件使用场景 \\
\midrule
CR0 & 保护模式、分页等控制位的一部分 & 决定地址和权限检查的总体模式 \\
CR3 & 页目录基址 & TLB 未命中或刷新后查页表 \\
GDTR & 全局描述符表位置和界限 & 装载段选择子、检查段属性 \\
IDTR & 中断/异常描述符表位置和界限 & 异常、中断、陷入入口选择 \\
EFLAGS & 条件码、方向、中断允许等状态 & 分支、算术结果和中断接收边界 \\
\bottomrule
\end{longtable}

\topic{从实模式进入保护模式是一段硬件约定切换}80386 复位后并不会自动处于完整分页保护环境。系统固件或早期启动代码必须先准备描述符表，再装载 GDTR，设置 CR0 中的保护模式相关位，并通过跳转刷新取指路径和段状态。分页还需要准备页目录、页表，把 CR3 指向页目录，再打开分页控制位。每一步都是硬件状态切换，顺序错误会导致取指、数据访问或异常入口无法解释。

\begin{longtable}{L{0.18\linewidth}L{0.4\linewidth}L{0.32\linewidth}}
\toprule
阶段 & 硬件动作 & 错误后果 \\
\midrule
准备 GDT/IDT & 在内存中放置段描述符和异常入口描述符 & 段装载或异常入口失败 \\
装载 GDTR/IDTR & CPU 记录描述符表基址和界限 & 后续选择子无法被正确解释 \\
设置 CR0 & 开启保护模式相关控制 & 旧的取指和段状态需要刷新 \\
远跳转/重装段 & 让 CS/DS/SS 等进入新描述符语义 & 继续按旧语义执行会失控 \\
准备页表并设置 CR3 & 给分页硬件页目录入口 & TLB 未命中时无法 walk \\
开启分页 & 线性地址开始经页表转换 & 映射缺失会立即异常 \\
\bottomrule
\end{longtable}

\topic{描述符和异常错误码把失败原因带进硬件入口}保护模式中的段描述符不只是“基址和长度”。它还包含类型、存在位、特权级、可读写/可执行等属性。分页项也不只是物理地址，还包含存在、读写、用户/特权、访问和脏等位。异常发生时，CPU 需要把足够信息交给异常入口，使系统能够判断是不存在、权限错误、越界，还是其他硬件条件失败。

\begin{longtable}{L{0.2\linewidth}L{0.34\linewidth}L{0.36\linewidth}}
\toprule
对象 & 关键字段 & 硬件问题 \\
\midrule
段描述符 & base、limit、type、DPL、present & 这个选择子是否存在、能否以当前特权使用 \\
页表项 & page frame、present、RW、US、accessed、dirty & 这个线性页是否存在、能否按本次方式访问 \\
异常入口 & selector、offset、gate type、DPL & 失败后 CPU 应进入哪段处理代码 \\
错误码 & 选择子或页错误原因位 & 处理程序如何定位失败对象和原因 \\
\bottomrule
\end{longtable}

一次页错误不是“内存读失败”这么简单。CPU 已经形成了线性地址，检查 PDE/PTE 或权限时发现不能继续，于是放弃普通写回路径，保存足够的异常信息，并从 IDT 中选择页错误入口。若处理器允许一部分状态已经提交、一部分状态未提交，系统就难以恢复。因此异常路径必须与指令提交边界配合：错误指令不应像成功指令一样修改目标寄存器或设备状态。

\topic{页错误诊断要同时看地址、访问类型和权限位}80386 把页错误从“总线读不到数据”提升为可诊断的硬件事件。CPU 会保存导致故障的线性地址，并通过异常错误码给出若干原因位。学习阶段不必记住所有后续扩展，但必须掌握三个基本问题：页是否存在，访问是读还是写，访问来自用户态还是特权态。只有把这三项与 PDE/PTE 中的 present、RW、US 等位对照，才能判断是缺页、写只读页、用户态访问内核页，还是页表本身没有建立。

\begin{longtable}{L{0.2\linewidth}L{0.34\linewidth}L{0.36\linewidth}}
\toprule
观察对象 & 含义 & 排查方向 \\
\midrule
故障线性地址 & 哪个线性地址无法完成转换或权限检查 & 页目录索引、页表索引和页内偏移是否落在预期区域 \\
P 位 & 故障由不存在页还是保护违规引起 & 缺页处理、页表初始化或权限设置 \\
W/R 位 & 本次访问是写还是读导致故障 & 写只读页、写保护代码段、copy-on-write 类机制 \\
U/S 位 & 本次访问来自用户级还是特权级 & 用户程序越权访问内核映射或段/页权限不一致 \\
IDT 入口 & 页错误应跳到哪个处理程序 & 异常入口描述符、栈切换和处理代码是否有效 \\
\bottomrule
\end{longtable}

这套诊断方法可以反过来指导最小保护模式实验。先建立一段恒等映射，让取指、栈和数据访问都可成功；再故意取消某个页表项的 present 位，确认访问该页会进入页错误入口；随后把某页设置为只读，执行写访问，确认错误码能区分权限失败和不存在页。若异常处理程序本身所在页没有映射，CPU 可能在处理一个错误时再次出错，最终表现为无法恢复的故障。因此保护模式 bring-up 必须先保证 GDT、IDT、栈、异常代码和页表所在内存可访问。

\topic{80386 总线的字节使能与 8086 一脉相承}80386 对外有 32-bit 数据总线。地址线通常给出 32-bit 对齐字的高位地址，字节使能信号指出本次传输中哪些 8-bit 字节有效。这样，一个字节写不会破坏同一 32-bit 数据片段中的其他字节；一个未对齐的 32-bit 写可能需要拆成多个传输，或由硬件和总线控制器处理额外周期。

\begin{longtable}{L{0.2\linewidth}L{0.28\linewidth}L{0.42\linewidth}}
\toprule
访问类型 & 有效字节车道 & 关键问题 \\
\midrule
字节访问 & 一个 BE 低有效 & 相邻字节不能被写改 \\
半字访问 & 两个相邻 BE 低有效 & 地址最低位和字节序必须一致 \\
对齐字访问 & 四个 BE 低有效 & 一个总线数据片段可完成传输 \\
未对齐字访问 & 跨越两个对齐片段 & 需要拆分周期或触发规定异常 \\
\bottomrule
\end{longtable}

从时延角度看，80386 到 80486/Pentium 的演进可以读成“把常用慢路径搬近”。TLB 避免每次访存都走页表；片上 cache 避免每次命中数据都走外部总线；流水线把长执行路径切成多个边界；分支预测减少控制流改变带来的等待。硬件并没有消灭延迟，而是用缓存、预测、并行和更细边界，让常见情况少走远路，让少见情况通过异常、回填或冲刷流水线处理。

从 8086 到 80386，可以看到 CPU 和整机边界的变化。8086 主要展示了早期微处理器如何通过外部总线接入存储器和 I/O；80386 则展示了处理器如何把内存管理和保护机制纳入芯片内部。后来的处理器继续把 cache、TLB、流水线、分支预测、乱序执行、多核互连等机制纳入芯片，但基本问题没有改变：地址从哪里来，权限如何检查，数据经哪条路径返回，状态在哪个边界提交。

80486 和 Pentium 可以作为 80386 之后的两个观察窗口。80486 把片上 cache、流水线和浮点方向的集成推到更显著的位置，使常用指令和数据更容易走近路径；Pentium 进一步展示双流水线、分支预测和更复杂控制对吞吐率的影响。它们并没有改变“寄存器到组合逻辑到寄存器”的基本模型，而是复制、切分和调度这些路径，让常见工作更少等待远处存储器和长控制路径。

\begin{longtable}{L{0.2\linewidth}L{0.32\linewidth}L{0.38\linewidth}}
\toprule
阶段 & 新增硬件重点 & 与前文概念的连接 \\
\midrule
80386 & 32-bit 地址、分页、保护和异常 & 地址路径和权限判断进入 CPU 内部关键结构 \\
80486 & 片上 cache、流水线、浮点集成 & 常见数据靠近核心，长组合路径被边界切分 \\
Pentium & 双流水线、预测和更复杂发射控制 & 多条候选执行路径需要统一调度和恢复 \\
现代核心 & 乱序、重命名、TLB、多级 cache、多核互连 & 快慢路径、状态提交和一致性成为核心问题 \\
\bottomrule
\end{longtable}

\section{三、整机启动与经典芯片后的演进}

通电或复位后，CPU、主存控制器、互连和设备控制器都必须进入规定初始状态。PC 通常被置到固定启动地址，那里连接 ROM、片上存储或其他启动介质。启动路径的第一件事，是让硬件从可预测位置开始取编码。

整机启动不仅是 CPU 复位。时钟源要稳定，复位树要按顺序释放，主存控制器要完成训练或初始化，互连要进入可转发请求的状态，必要设备要处于安全默认值。若这些硬件条件没有满足，即使 CPU 本身能取指，也可能在第一次访问主存或设备时无法继续推进。

\topic{从复位向量走到第一条有效指令}启动可以写成一条硬件 trace。上电后，电源管理电路等待电压进入允许范围；时钟源或晶振稳定后，复位电路保持 CPU 和关键外设在安全状态；复位释放时，PC 或取指地址逻辑装入固定复位向量；地址译码把这个地址指向 ROM、Flash 或片上启动存储；取指路径读回第一条编码；控制器译码并开始执行初始化序列。任何一步没有定义，整机都会表现为“CPU 不启动”，但根因可能不在 CPU。

\begin{longtable}{L{0.16\linewidth}L{0.42\linewidth}L{0.32\linewidth}}
\toprule
启动阶段 & 硬件动作 & 预期行为 \\
\midrule
电源爬升 & 电压达到允许范围，去耦提供瞬态电流 & 欠压复位未提前释放 \\
时钟稳定 & 晶振、PLL 或外部时钟进入稳定范围 & 时钟频率和占空比满足规格 \\
复位保持 & CPU、控制器、设备进入安全默认状态 & 危险输出关闭，写使能无效 \\
复位释放 & PC 装入复位向量或启动地址 & 第一拍地址可预测 \\
取启动码 & ROM/Flash/片上存储响应该地址 & 数据总线返回有效编码 \\
初始化 & 配置栈、RAM、时钟、设备和中断 & 每项配置有完成或错误状态 \\
\bottomrule
\end{longtable}

\begin{pdfdiagram}
  \node[hw state, minimum width=1.45cm] (pwr) at (0,0) {电源};
  \node[hw compute, minimum width=1.55cm] (clk) at (1.95,0) {时钟};
  \node[hw control, minimum width=1.65cm] (rst) at (4.0,0) {复位树};
  \node[hw state, minimum width=1.7cm] (vec) at (6.15,0) {复位向量};
  \node[hw compute, minimum width=1.75cm] (rom) at (8.35,0) {Boot ROM};
  \node[hw state, minimum width=1.55cm] (fw) at (10.45,0) {固件};
  \draw[hw bus] (pwr) -- node[above=2pt,hw label,fill=white,inner sep=1pt]{稳定} (clk);
  \draw[hw bus] (clk) -- node[above=2pt,hw label,fill=white,inner sep=1pt]{可用} (rst);
  \draw[hw bus] (rst) -- node[above=2pt,hw label,fill=white,inner sep=1pt]{释放} (vec);
  \draw[hw bus] (vec) -- node[above=2pt,hw label,fill=white,inner sep=1pt]{取指地址} (rom);
  \draw[hw bus] (rom) -- node[above=2pt,hw label,fill=white,inner sep=1pt]{第一条编码} (fw);
  \node[hw label, text width=9.5cm] at (5.2,-1.05) {启动不是“CPU 开始跑”一句话，而是一条逐级成立的链路：电源和时钟先可靠，复位再释放，复位向量必须命中非易失启动代码。};
\end{pdfdiagram}
\diagramnote{整机启动链路。第一条有效指令依赖电源、时钟、复位、地址映射和启动存储共同成立。}

\topic{复位入口体现地址空间顶端和底端的不同设计}不同处理器的复位入口不相同，但背后的硬件问题一致：复位后 CPU 必须从非易失存储中的确定位置取得第一条可执行编码。6502 从固定向量中取入口地址，Z80/8051 常从 0 地址开始执行，8086 从 1 MiB 空间顶端附近开始取指，80386 及后续 x86 则把初始取指放到更高物理地址附近，以便系统固件映射在地址空间顶端。

\begin{longtable}{L{0.2\linewidth}L{0.28\linewidth}L{0.42\linewidth}}
\toprule
处理器 & 典型复位入口 & 硬件含义 \\
\midrule
6502 & 从 \code{0xFFFC/0xFFFD} 取 16-bit 入口 & 顶端向量表必须由 ROM 响应 \\
Z80/8051 & 从 \code{0x0000} 开始执行 & 低地址必须映射启动代码 \\
8086 & 物理地址 \code{0xFFFF0} 附近 & 1 MiB 顶端 ROM 保存跳转或启动片段 \\
80386 及后续 x86 & 高物理地址顶端附近 & 固件映射到高地址，早期代码再跳转到初始化区域 \\
现代 SoC & 片上 Boot ROM 或可配置启动介质 & 先验证/选择介质，再初始化 DRAM 和外设 \\
\bottomrule
\end{longtable}

在 8086 类系统中，复位后从规定的高地址入口开始取指，外部 ROM 必须映射到这个入口；在许多微控制器中，复位向量可能位于片内 Flash 的固定地址，前几个字保存初始栈指针和入口地址；在现代 SoC 中，片上 Boot ROM 可能先验证启动介质、配置内存控制器，再把控制权交给后续固件。形式不同，问题相同：复位向量指向哪里，谁在该地址响应，第一批指令执行前哪些硬件已经可靠。

\topic{整机实验可以从 ROM、RAM 和两个 MMIO 寄存器开始}一个可搭建的最小整机不需要复杂操作系统。可先规定 8-bit 或 16-bit 地址空间：低地址接 ROM，中间接 RAM，高地址映射 LED 输出寄存器和按键输入寄存器。CPU 或教学控制器复位后从 ROM 取指，执行一段程序：读取按键状态，若按下则把 LED 寄存器置 1，否则置 0。这个系统已经包含整机硬件的核心对象：取指、数据访存、地址译码、MMIO、副作用、外部输入同步和可观察输出。

\begin{lstlisting}
minimal whole-machine program:
  loop:
    r1 = MMIO[GPIO_IN]
    if r1 == 0 goto off
    MMIO[GPIO_OUT] = 1
    goto loop
  off:
    MMIO[GPIO_OUT] = 0
    goto loop
\end{lstlisting}

这段程序的硬件检查不应只看 LED 是否跟随按钮。应同时确认 ROM 地址、RAM 地址和 MMIO 地址不会被同时片选；按钮已经同步后才被读；写 \code{GPIO_OUT} 只改变输出锁存器，不误写 RAM；分支改变的是 PC 下一值；循环运行时没有总线争用。若以后把 LED 换成电机、继电器或功率负载，还必须增加驱动级、续流保护和安全停机链路，不能把 MMIO 输出锁存器直接当成功率开关。

\topic{最小整机的检查应分成三层}第一层检查硬件静态条件：电源、复位、时钟、片选互斥和三态方向。第二层检查单条指令：ADD、MOV、JZ 等每条指令都能在单步下看到正确的控制线和提交边界。第三层检查小程序：ROM 中循环程序能够读按键、分支、写 LED，并能在异常输入下保持安全状态。

\begin{longtable}{L{0.18\linewidth}L{0.42\linewidth}L{0.3\linewidth}}
\toprule
层次 & 具体检查 & 失败时优先排查 \\
\midrule
静态硬件 & 电源纹波、复位释放、时钟边沿、片选互斥 & 供电、复位树、译码极性 \\
单指令 & PC、IR、控制字、ALUOut、MDR、写回寄存器 & 控制器、寄存器堆、ALU、存储器时序 \\
小程序 & 按键输入同步、条件分支、LED 输出锁存 & MMIO 地址、分支目标、外设副作用 \\
错误路径 & 未映射地址、慢设备等待、外设 error 位 & 超时、错误响应、异常入口 \\
\bottomrule
\end{longtable}

在实验报告中，应记录每个检查点的“预期值、实测值、观察工具和结论”。例如 PC 第 0 拍应为 \code{0x0000}，ROM 片选应有效，IR 应锁存第一条 \code{MOV}；执行写 \code{GPIO_OUT} 的指令时，RAM 片选应无效，LED 锁存器写使能应只出现一个有效窗口。这样的记录比“程序运行成功”更接近工程实践。

\topic{可采购器件清单要匹配教学目标}若目标是搭一台低速可观察整机，器件选择应服务于可测性，而不是追求最高频率。可以使用 74HC/74HCT 系列计数器、锁存器、三态缓冲器、加法器、比较器、EEPROM、SRAM 和 LED/按键模块；也可以用 CPLD/FPGA 把控制器和数据通路放进可编程逻辑，再把 ROM/RAM/MMIO 作为独立模块观察。

\begin{longtable}{L{0.2\linewidth}L{0.34\linewidth}L{0.36\linewidth}}
\toprule
对象 & 可选器件 & 选择理由 \\
\midrule
PC/计数 & 74HC161/163 或 FPGA 寄存器 & 可清零、可装载、可单步观察 \\
IR/锁存 & 74HC574/273 & 时钟沿保存指令或中间状态 \\
ALU 原型 & 74HC283、异或门、比较器或 FPGA 逻辑 & 可先实现加法、零检测和简单逻辑 \\
总线隔离 & 74HC245/244 & 控制数据方向，避免多源驱动 \\
ROM & AT28C 系列 EEPROM 或片内 ROM & 保存启动程序和指令编码 \\
RAM & 低速异步 SRAM 或 FPGA block RAM & 接口直观，便于单步读写 \\
I/O & LED 锁存器、按键同步器、UART 模块 & 形成可观察输入输出和串口调试 \\
\bottomrule
\end{longtable}

搭建顺序应从电源、时钟、复位、单个寄存器和单条总线开始。只有确认三态控制不会冲突、复位后所有写使能无效、时钟边沿干净，才应接入 ROM、RAM 和 I/O。若直接把所有模块连接后调试，任何错误都可能表现为“程序不跑”，排查成本会急剧上升。

\topic{整机 bring-up 要记录症状、信号和假设}真实硬件第一次上电时，常见问题不是“CPU 坏了”，而是复位、时钟、片选、总线方向、等待状态或启动 ROM 映射有误。把症状、相关信号和当前假设逐项记录下来，才能形成一条可逐步核对的排查链条，而不是笼统地说“检查硬件”。

\begin{longtable}{L{0.2\linewidth}L{0.34\linewidth}L{0.36\linewidth}}
\toprule
症状 & 优先观察信号 & 可能原因 \\
\midrule
PC 不变化 & 复位、时钟、PC 写使能 & 复位未释放、时钟不稳定、控制器未进入取指 \\
ROM 无输出 & 地址线、ROM 片选、输出使能 & 复位向量未映射到 ROM，译码极性错误 \\
数据总线异常发热 & CPU/存储器输出使能、方向控制 & 两个器件同时驱动总线 \\
加载读错值 & 地址、字节使能、读响应、MDR & 地址形成错误或存储器时序不满足 \\
LED 乱闪 & MMIO 片选、写使能、外设锁存时钟 & 地址译码过宽或写周期毛刺 \\
中断重复进入 & 设备状态、中断控制器 ACK、CPU 入口 & 清除顺序错误或事件源未确认 \\
\bottomrule
\end{longtable}

经典微处理器的发展可以用若干硬件剖面概括。下表不是处理器史年表，而是说明每一代芯片把哪些硬件问题推到读者面前。

\begin{longtable}{L{0.18\linewidth}L{0.2\linewidth}L{0.5\linewidth}}
\toprule
芯片或系列 & 典型特征 & 对硬件学习的意义 \\
\midrule
4004 & 4-bit 微处理器 & 微处理器把控制器、寄存器和算术部件集成到单芯片形态 \\
6502/Z80 & 经典 8-bit 微处理器 & 外部总线、存储器、I/O 和多周期控制构成早期微机基本形态 \\
8051 & 经典单片机 & CPU、片内存储、I/O、定时器和中断控制展示片上整机控制 \\
8080/8085 & 8-bit 微处理器 & 外部总线、地址空间、存储器和 I/O 开始形成通用微机结构 \\
8086/8088 & 16-bit x86 起点 & 段:偏移地址、指令预取、复用总线和外部锁存器体现早期整机取舍 \\
80286 & 保护模式前驱 & 分段保护和特权机制开始成为处理器硬件的一部分 \\
80386 & 32-bit x86 & 32-bit 寄存器、线性地址、分页和异常机制把 CPU 推向系统级硬件 \\
80486/Pentium & 片上 cache 与更深执行结构 & 流水线、cache、浮点单元和更复杂控制逻辑成为性能核心 \\
\bottomrule
\end{longtable}

8051 的意义在于把“小整机”收进一颗芯片。它让读者看到，CPU 旁边可以直接布置片内 RAM、ROM、I/O 口、定时器、串口和中断控制。外部引脚不再只是地址/数据总线，也可能直接是可编程 I/O。对硬件学习而言，这说明整机边界可以移动：同样是“写设备寄存器改变输出”，在早期微机上可能通过板级地址译码访问外设，在单片机上则可能是写特殊功能寄存器改变某个引脚锁存器。

80486 的意义在于把片上 cache 和更清楚的流水化执行推到学习视野中。片上 cache 缩短了高频取指和数据访问路径，使常见访问不必每次走外部总线；流水线把取指、译码、执行等工作分段，使不同指令可以重叠推进。它同时提醒读者，性能不是单靠提高门速度，\hwevidence{把常见数据放近}和\hwtiming{把长路径切段}同样关键。

Pentium 的意义在于展示更强的并行执行方向。双流水线、分支预测和更复杂的发射控制，使处理器尝试在同一时间推进多条合适的指令。此时硬件问题不再只是“一个 ALU 算得多快”，还包括指令之间是否有依赖、分支方向是否预测正确、两条流水线是否争用资源、错误路径如何清除、异常如何保持精确边界。这些机制为现代乱序执行、寄存器重命名和提交队列铺路。

\topic{80486 把 cache、FPU 和流水线收进芯片}80386 时代，cache 是主板上的独立芯片加标签存储器，浮点运算是另一颗 80387 协处理器，两者都通过板级总线与 CPU 通信。80486 把三样东西收进片内：8 KiB 统一 cache（指令和数据共用，采用写穿策略）、浮点单元和一条 5 级流水线。配套的还有总线接口的改变：486 引入突发总线周期，一次 cache 行填充用 4 拍连续传完 16 B，只有第一拍需要给出地址，后面三拍不再重复。收益很直接：cache 命中时取指和读数不再走板级总线，常见指令的吞吐接近每拍一条；代价同样明确：片内 cache 要处理一致性，写穿策略持续占用外部总线带宽，流水线的段间控制、数据相关和分支处理也比多周期实现复杂得多。

\topic{Pentium 用双发射和分支预测重组执行结构}Pentium 的性能提升不再主要来自时钟频率，而是来自执行结构的重组。它有 U、V 两条整数流水线，满足配对规则的简单指令可以在同一拍内双发射；分支目标缓冲（BTB）预测分支方向，猜错才冲刷流水线；指令 cache 和数据 cache 分开，各 8 KiB，取指和读数不再争同一个端口；外部数据总线扩到 64 bit，一次总线事务搬运的数据翻倍。Pentium Pro 又往前走了一步：x86 指令先被译成微操作，由乱序执行核心按数据就绪顺序执行，寄存器重命名消除假依赖，最后按程序顺序提交，从而在乱序执行下保持精确异常。执行可以和提交分离，是这一代留下的主线：第四章讲的执行/提交区分，在这里成为真实芯片的基本组织方式。

\topic{386/486 主板把 CPU、内存和外设分成平台层次}80386、80486 之后，整机越来越依赖主板芯片组。CPU 不再直接以简单板级总线面对所有外设，而是通过内存控制器、cache 控制器、总线桥和 I/O 控制器访问不同速度、不同协议的目标。早期 PC 主板可以粗略理解为：CPU 与高速本地总线连接，芯片组负责 DRAM、cache、ISA/EISA/VLB/PCI 等扩展路径，低速外设再挂到更慢的 I/O 控制器。后来的“北桥/南桥”说法，就是把高速内存/图形路径和低速 I/O 路径分开组织。

\begin{longtable}{L{0.2\linewidth}L{0.34\linewidth}L{0.36\linewidth}}
\toprule
平台层次 & 负责对象 & 与本章概念的关系 \\
\midrule
CPU 本地总线 & 取指、数据访问、cache miss、写回 & 地址、数据、字节使能、等待和总线错误 \\
内存/cache 控制 & DRAM 行列访问、二级 cache、写缓冲 & 命中/未命中、回填、写回和时序预算 \\
高速扩展路径 & 图形、磁盘控制、网络或总线桥 & DMA、总线主设备、仲裁和中断 \\
低速 I/O 控制 & 串口、并口、键盘、软驱、RTC 等 & 端口 I/O、状态寄存器、完成事件 \\
固件层 & BIOS/Option ROM、设备初始化和启动选择 & 复位向量、枚举、内存检测和启动介质 \\
\bottomrule
\end{longtable}

这个分层解释了为什么整机性能不能只看 CPU 主频。若 cache 未命中要走慢 DRAM，若 PCI 设备 DMA 被总线仲裁阻塞，若中断控制器或桥片配置错误，CPU 再快也会等待。另一方面，把 cache、内存控制器、图形控制、PCIe 根复合体等功能逐步移近 CPU，也正是现代平台降低时延的方向：常用路径减少芯片间往返，慢速外设通过桥接和队列异步处理。

\topic{从 PCI 到 PCIe，设备仍然是事务发起者和响应者}PCI/PCIe 这类现代扩展互连比早期 ISA 总线复杂得多，但本章的分析方法仍然适用。设备可以暴露配置空间、MMIO BAR、DMA 能力和中断机制；CPU 或固件先枚举设备，分配地址窗口和中断资源；驱动程序再通过 MMIO 配置设备，通过 DMA 描述符让设备读写内存，并通过 MSI/MSI-X 或传统中断接收完成事件。

\begin{longtable}{L{0.18\linewidth}L{0.38\linewidth}L{0.34\linewidth}}
\toprule
PCI/PCIe 对象 & 硬件含义 & 关键问题 \\
\midrule
配置空间 & 设备身份、能力、BAR 和控制位 & 枚举是否能读到设备，BAR 是否分配到不冲突地址 \\
BAR/MMIO & 设备寄存器窗口映射到物理地址 & 访问属性是否不可缓存或按设备规则排序 \\
DMA & 设备作为总线发起者读写主存 & 地址权限、IOMMU、cache 一致性和 completion 顺序 \\
MSI/MSI-X & 设备通过内存写形式投递中断消息 & 消息地址/数据配置是否正确，是否会丢事件 \\
错误报告 & 链路、权限、超时或设备内部错误 & 错误能否被记录、上报和恢复 \\
\bottomrule
\end{longtable}

这一段对学习现代硬件很重要：PCIe 不再是多块板卡共享一排并行地址线的简单总线，但它仍然要解决同样的问题。请求从哪里来，目标地址是谁，数据何时有效，完成如何返回，错误如何上报，中断如何投递，DMA 写入何时对 CPU 可见。只要把这些问题标出来，复杂协议就不会变成无法理解的名词堆。

\topic{SoC 把北桥、南桥和许多外设继续收进芯片}现代 SoC 的趋势是把 CPU 核、cache、内存控制器、GPU/NPU、视频编解码、PCIe/USB/SATA、网卡、显示控制、GPIO、定时器和安全启动逻辑放进一颗或少数几颗芯片。这样做的直接收益是缩短常见路径、降低引脚数量、减少板级功耗和提高集成度；代价是片上互连、一致性、时钟/电源域、调试和安全启动变得更复杂。

\begin{longtable}{L{0.2\linewidth}L{0.34\linewidth}L{0.36\linewidth}}
\toprule
SoC 结构 & 低时延来源 & 新增边界 \\
\midrule
片上 L2/L3 cache & 多个核心共享近处数据 & 一致性状态、替换、QoS 和带宽竞争 \\
集成内存控制器 & CPU 到 DRAM 控制路径更短 & DRAM 训练、刷新、ECC 和功耗状态 \\
片上互连 & 核、DMA、GPU、外设通过网络通信 & 仲裁、优先级、背压和死锁避免 \\
集成外设 & GPIO、UART、SPI、I2C、USB 等近距离访问 & MMIO 属性、中断路由和复位/时钟门控 \\
安全启动 & Boot ROM 验证后续固件 & 密钥、签名、回滚保护和调试锁定 \\
\bottomrule
\end{longtable}

因此，“现代芯片为什么时延低”可以更完整地回答：不是所有路径都低时延，而是常见路径被集成到片上、被 cache/TLB 缓存、被预取和预测提前准备、被流水线和乱序执行并行化、被 DMA 和队列卸载；不常见路径则通过 cache miss、TLB miss、DRAM 排队、设备背压、错误报告和中断恢复来处理。低时延来自体系结构、微架构、互连、存储层级和平台集成共同作用。

从 8086、80386 向后看，处理器继续沿几条方向扩展。第一，内部执行路径更深，流水线把取指、译码、执行、访存、写回切成多个阶段。第二，存储层级更复杂，片上 cache 和 TLB 成为地址访问性能的关键。第三，控制逻辑更复杂，分支预测、乱序执行和异常恢复要求硬件保存更多内部状态。第四，多核和片上互连让“总线事务”扩展成核间一致性和缓存一致性协议。第五，外设从简单寄存器发展到 PCI、PCIe、USB、SATA 等分层协议，但本质仍然是地址、数据、控制、握手、完成和错误状态。

\begin{longtable}{L{0.24\linewidth}L{0.58\linewidth}}
\toprule
演进方向 & 保持不变的硬件问题 \\
\midrule
流水线 & 状态边界如何切分，错误路径如何清除 \\
cache/TLB & 地址如何命中、缺失、替换和回填 \\
保护与异常 & 何时拒绝访问，如何进入规定处理入口 \\
DMA 与设备 & 谁发起事务，谁完成事务，完成如何通知 CPU \\
多核一致性 & 多个观察者如何看到受控顺序的内存结果 \\
\bottomrule
\end{longtable}

到这里，本书的硬件链路闭合：电和器件形成可控开关，开关形成门，门形成组合逻辑，反馈和时钟形成状态，寄存器和 ALU 形成数据通路，控制器组织动作，CPU 执行编码，主存、cache、互连、设备和启动路径把它扩展成整机。8086 和 80386 的意义在于提供了两个清晰历史剖面：一个展示早期微处理器如何连接总线和地址空间，一个展示处理器如何承担更复杂的地址转换、保护和异常机制。

\section{四、从目标文件、固件到可启动机器}

前面几章已经解释了 CPU 如何执行编码，主存和设备如何响应地址，经典芯片如何定义平台边界。还差一条常被硬件教材略过、却对系统理解非常关键的链路：程序怎样从目标文件进入内存，符号地址怎样变成机器地址，启动固件怎样把硬件带到可执行状态，异常表和页表怎样让失败路径可恢复。参考经典系统教材的读法，这一节不展开编译器细节，而是把链接、装载和启动都还原成硬件地址约定。

\topic{目标文件里的符号最终要变成地址}
汇编器或编译器产生的目标文件通常还不知道所有最终地址。函数名、全局变量名、外部符号和跳转目标先以符号或重定位项表示。链接器把多个目标文件和库合并，决定代码段、只读数据段、数据段、BSS 段等放到什么地址，并把需要修补的位置改成最终地址或相对偏移。

\begin{longtable}{L{0.18\linewidth}L{0.38\linewidth}L{0.34\linewidth}}
\toprule
对象 & 系统含义 & 硬件落点 \\
\midrule
代码段 & 机器指令编码 & 取指路径从这些地址读取 opcode \\
只读数据 & 常量、跳转表、字符串 & 普通加载读取，权限通常禁止写 \\
数据段 & 已初始化全局变量 & 启动时从镜像复制到 RAM \\
BSS & 零初始化全局变量 & 启动时清零，镜像中不必保存全量零字节 \\
重定位项 & 需要链接器或装载器修补的地址位置 & 改写指令立即数字段、地址表或指针字 \\
\bottomrule
\end{longtable}

这说明“函数地址”不是抽象标签。取指硬件需要 PC 中的数值地址；CALL 需要目标地址或相对偏移；全局变量加载需要数据地址；中断向量需要入口地址。链接器和装载器的工作，就是把这些名字和段布局变成 CPU 能取指、加载、存储和跳转的地址。

\begin{lstlisting}
source idea:
  call uart_putc
  load global_counter

after link/load:
  CALL 0x00102480
  MOV R1, [0x00203010]
\end{lstlisting}

\topic{装载器建立的是第一张可执行地址地图}
在有操作系统的机器上，装载器把可执行文件映射进进程地址空间，建立代码、数据、堆、栈和共享库区域；在裸机或固件场景中，启动代码负责把镜像从 ROM/Flash 复制或映射到合适位置，清 BSS，设置栈指针，初始化全局数据，然后跳到入口函数。两者规模不同，但硬件问题一致：哪些地址可取指，哪些地址可读写，栈在哪里，异常入口在哪里，设备寄存器在哪里。

\begin{longtable}{L{0.18\linewidth}L{0.42\linewidth}L{0.3\linewidth}}
\toprule
启动/装载动作 & 硬件意义 & 失败表现 \\
\midrule
设置栈指针 & CALL/RET、局部变量和异常入口有可用内存 & 第一次函数调用或中断即破坏内存 \\
复制数据段 & RAM 中全局变量得到初始值 & 程序读到未初始化数据 \\
清 BSS & 零初始化对象符合语言和固件约定 & 状态机、计数器、指针初值随机 \\
建立向量表 & 中断和异常有入口地址 & 外部事件或错误无法处理 \\
初始化时钟/总线 & 后续 RAM、设备、cache 路径满足时序 & 访问慢器件超时或读写失败 \\
跳入口 & PC 改到主程序第一条指令 & 程序真正进入可维护代码路径 \\
\bottomrule
\end{longtable}

裸机启动常见伪代码如下。它表面像软件初始化，实质上每一步都在建立后续硬件事务的前提。

\begin{lstlisting}
reset_handler:
  set SP to top_of_stack
  copy .data from flash image to SRAM
  zero .bss in SRAM
  initialize clocks and memory controller
  install vector table base
  initialize essential MMIO devices
  jump main
\end{lstlisting}

\topic{异常向量表是硬件失败路径的地址表}
只要系统会处理中断、页错误、非法指令或设备事件，就需要某种入口表。早期处理器可能用固定地址或固定向量；保护模式处理器使用更复杂的描述符表；微控制器常在启动区域放置向量表。表的形式不同，但作用相同：当硬件事件发生时，CPU 能把事件编号转成处理程序入口，并保存足够状态以便恢复或诊断。

\begin{longtable}{L{0.18\linewidth}L{0.36\linewidth}L{0.36\linewidth}}
\toprule
向量项 & 保存内容 & 硬件使用场景 \\
\midrule
复位向量 & 第一条启动代码地址或跳转指令 & 上电或复位后装入 PC \\
非法指令向量 & 无法译码或不允许执行时的入口 & opcode 不合法或特权不允许 \\
页错误/总线错误向量 & 地址转换或外部事务失败入口 & 加载/存储/取指不能完成 \\
外设中断向量 & 设备完成、错误或数据到达入口 & 中断控制器投递事件 \\
系统调用入口 & 用户代码请求特权服务入口 & 权限切换和受控进入内核或监控程序 \\
\bottomrule
\end{longtable}

向量表本身也必须位于可访问地址。若页表或地址译码没有覆盖异常入口，CPU 在处理第一个错误时会遇到第二个错误；若栈没有准备好，保存 PC/Flags 或错误码时会破坏未知内存；若中断处理程序所在代码被错误标成不可执行，事件到达后系统会在入口处再次失败。启动顺序必须先保证最小异常路径可靠，再打开复杂设备和中断。

\topic{页表把装载结果提升为可保护的运行时地图}
链接器决定镜像中的相对布局，装载器把段放进内存，页表则把运行时地址空间变成可检查的约定。代码页可以标为只读和可执行，数据页可读写但不可执行，内核页禁止用户访问，设备 MMIO 页设置为不可缓存或强顺序。这样，同一条加载/存储/取指路径在真正访问 cache 或总线前，会先经过地址转换和权限检查。

\begin{longtable}{L{0.2\linewidth}L{0.36\linewidth}L{0.34\linewidth}}
\toprule
区域 & 典型页属性 & 硬件效果 \\
\midrule
代码 & 只读、可执行 & 对代码页的写操作触发保护异常 \\
只读数据 & 只读、不可执行或按平台规定 & 防止常量被误写或执行数据 \\
普通数据/堆/栈 & 可读写、通常不可执行 & 数据写入允许，取指可被阻止 \\
内核映射 & 特权访问，用户禁止 & 用户态越权访问触发异常 \\
MMIO 窗口 & 不可缓存、设备顺序属性 & 读写按设备事务执行，不被普通 cache 吞掉 \\
\bottomrule
\end{longtable}

这张表把前几章内容合并起来：第三章的 bit 和地址解释，第四章的 PC 与加载/存储， 第五章的 cache/MMIO 属性，第六章的 80386 分页和异常，在页表属性中汇合。现代系统的许多安全和可靠性机制，本质上就是让硬件在每次取指和访存前检查这些约定。

\topic{从固件到操作系统是平台所有权转移}
固件启动阶段拥有硬件控制权：它配置时钟、内存控制器、基本 I/O、启动介质和必要安全检查。操作系统接管后，会重新建立页表、中断控制、设备驱动和调度机制。这个交接不是一句“加载内核”，而是一段平台所有权转移：谁拥有中断表，谁拥有页表，谁管理内存，谁接管设备，谁定义错误处理策略。

\begin{longtable}{L{0.18\linewidth}L{0.4\linewidth}L{0.32\linewidth}}
\toprule
阶段 & 固件职责 & 交给后续系统的约定 \\
\midrule
早期启动 & 电源、时钟、复位、最小取指路径 & CPU 能从可信入口执行 \\
内存初始化 & DRAM 控制器、cache/TLB 初始策略 & 可用内存范围和保留区域 \\
设备发现 & 基本总线枚举、启动介质选择 & 设备地址窗口、中断和 DMA 能力 \\
镜像验证/加载 & 读取、校验、解压或复制内核/程序 & 入口地址、参数和内存布局 \\
控制权转移 & 设置寄存器、栈、页表或模式后跳转 & 后续系统接管异常、中断和设备 \\
\bottomrule
\end{longtable}

在微控制器上，固件和“应用”可能就是同一个镜像；在 PC 或服务器上，固件、引导加载器、内核和用户程序分层明显；在 SoC 上，还可能有 Boot ROM、一级引导、可信固件、安全监控和普通操作系统。层数不同，但每一层都要把可执行地址、栈、异常入口、内存属性和设备所有权交代清楚。

\topic{安全启动把第一条指令也纳入信任链}
现代 SoC 常把 Boot ROM 固化在芯片内部，复位后先执行不可修改的 ROM 代码。Boot ROM 读取下一阶段固件，验证签名、版本和配置，再决定是否跳转。这样做不是密码学装饰，而是硬件启动约定的一部分：如果第一段可变固件不可信，后续页表、设备配置、密钥和操作系统都可能被篡改。

\begin{longtable}{L{0.18\linewidth}L{0.4\linewidth}L{0.32\linewidth}}
\toprule
安全启动对象 & 硬件/固件动作 & 失败处理 \\
\midrule
Boot ROM & 从固定复位入口执行，不可现场改写 & 提供根信任和最小验证代码 \\
公钥/哈希 & 存在熔丝、ROM 或安全存储中 & 用于验证下一阶段镜像 \\
镜像签名 & 覆盖代码、配置和版本信息 & 验证失败则停止、降级或进入恢复模式 \\
回滚保护 & 记录允许的最小版本 & 防止加载旧漏洞固件 \\
调试锁定 & 限制 JTAG、串口下载或内存读出 & 防止启动后被外部接口改写控制流 \\
\bottomrule
\end{longtable}

这也解释了为什么“硬件从零到整机”不能只讲电路能跑。整机一旦能取指、访存和访问设备，就必须回答谁能改代码、谁能访问内存、谁能发起 DMA、谁能接收中断、启动镜像是否可信、错误路径是否可恢复。经典芯片提供了这些问题的历史剖面，现代平台则把它们推向更复杂的安全和系统集成边界。

\topic{案例：一个裸机镜像从链接地址跑到 main}
假设一块教学板把 Flash 映射到 \code{0x00000000}，SRAM 映射到 \code{0x20000000}，GPIO 映射到 \code{0x40000000}。链接脚本规定向量表和代码放在 Flash，已初始化数据运行在 SRAM，但初始内容存放在 Flash 镜像中，BSS 在 SRAM 中清零，栈顶放在 SRAM 高地址。这个布局看起来是软件文件，实际上直接决定复位后每一次取指、加载和存储的地址。

\begin{lstlisting}
memory map:
  0x00000000-0x0003FFFF  Flash: vectors, code, const, data image
  0x20000000-0x20007FFF  SRAM: data, bss, stack
  0x40000000-0x40000FFF  GPIO MMIO
\end{lstlisting}

\begin{longtable}{L{0.18\linewidth}L{0.38\linewidth}L{0.34\linewidth}}
\toprule
镜像区域 & 链接/装载含义 & 硬件事务 \\
\midrule
向量表 & 复位入口和异常入口地址 & 复位时取入口，异常时取向量 \\
\code{.text} & 指令编码放在 Flash & PC 取指访问 Flash 或映射后的代码区 \\
\code{.rodata} & 常量和只读表 & 加载读取，通常禁止写 \\
\code{.data} 镜像 & 初始值存放在 Flash & 启动代码复制到 SRAM 运行地址 \\
\code{.bss} & 运行时应为 0 的对象 & 启动代码对 SRAM 执行写操作清零 \\
栈 & 函数调用和异常保存上下文 & SP 指向 SRAM 可写区域 \\
\bottomrule
\end{longtable}

复位后第一段代码不应依赖尚未初始化的全局变量，也不应调用需要完整栈和运行时的复杂函数。它要先建立最小可运行边界。

\begin{lstlisting}
reset trace:
  PC = vector_table.reset_handler
  SP = _stack_top
  copy bytes from _sidata in Flash to _sdata in SRAM
  zero bytes from _sbss to _ebss
  configure clocks and essential wait states
  set vector table base if architecture supports it
  call main
\end{lstlisting}

\begin{longtable}{L{0.16\linewidth}L{0.42\linewidth}L{0.32\linewidth}}
\toprule
启动动作 & 不能省略的原因 & 常见错误 \\
\midrule
设置 SP & 后续 CALL、PUSH、异常保存需要栈 & SP 指到 Flash 或未映射区 \\
复制 \code{.data} & 全局初值要出现在 RAM 运行地址 & 变量初值全错或仍为 Flash 内容 \\
清 \code{.bss} & 语言和固件约定零初始化 & 状态机初值随机，指针非零 \\
配置等待状态 & Flash、SRAM 或外设满足时序 & 高频下取指或访问偶发失败 \\
设置向量表 & 中断和异常能进入正确入口 & 第一场中断跳到错误地址 \\
跳 \code{main} & 进入应用主逻辑 & 入口地址或模式错误导致程序失控 \\
\bottomrule
\end{longtable}

\topic{案例：保护模式最小启动清单}
以 80386 风格机器为例，若要从实模式进入带分页的保护环境，必须先建立一组硬件可读的表和寄存器状态。这个过程常被写成启动代码技巧，但本质是把 CPU 的地址解释约定从旧模式切换到新模式。

\begin{longtable}{L{0.16\linewidth}L{0.42\linewidth}L{0.32\linewidth}}
\toprule
步骤 & 硬件目的 & 预期行为 \\
\midrule
建立 GDT & 给代码段、数据段、栈段描述符 & 描述符 base/limit/type/present 正确 \\
建立 IDT & 给异常和中断入口描述符 & 页错误、非法指令等入口可达 \\
加载 GDTR/IDTR & CPU 知道描述符表位置 & 表所在内存可读，界限正确 \\
打开保护模式位 & 改变段选择子解释方式 & 远跳转后 CS 使用新描述符 \\
重装段寄存器 & DS/SS/ES 等进入新约定 & 数据和栈访问按新描述符解释 \\
建立页目录/页表 & 线性地址可转换到物理地址 & 取指、栈、异常代码均被映射 \\
加载 CR3 并开启分页 & CPU 开始查页表和 TLB & 缺失映射会触发页错误而非随机访问 \\
\bottomrule
\end{longtable}

最小保护模式 bring-up 应先保守映射。把当前代码、栈、页表、GDT、IDT 和 VGA/串口调试 MMIO 映射好，再打开分页；确认能打印或点亮调试输出后，再逐步收紧权限。若一开始就建立复杂高半区映射、多个页大小、用户/内核分离和设备属性，任何一个地址错误都可能导致连续异常，难以定位。

\begin{lstlisting}
protected-mode checklist:
  GDT covers code/data/stack selectors
  IDT covers at least fault handlers
  stack is writable after mode switch
  current instruction stream is mapped
  page tables themselves are mapped
  debug output MMIO is mapped uncached
  fault handlers do not fault recursively
\end{lstlisting}

\topic{案例：从异常号定位启动失败}
启动阶段最可怕的现象是“黑屏”或“无输出”。更有效的做法是把失败分层。若能进入非法指令异常，说明取指至少到了某条不合法编码；若能进入页错误，说明保护模式和 IDT 至少部分工作；若连续进入双重错误或复位，说明异常处理路径本身可能失败；若完全没有总线活动，可能是复位、时钟或复位向量映射问题。

\begin{longtable}{L{0.22\linewidth}L{0.38\linewidth}L{0.3\linewidth}}
\toprule
现象 & 可能边界 & 优先观察 \\
\midrule
无取指总线活动 & 复位、时钟、复位向量、ROM 片选 & 示波器看时钟/复位/地址线 \\
取指后立即异常 & 编码、模式、段描述符或权限错误 & IR/PC、异常号、错误码 \\
页错误循环 & 页表缺映射、权限错误、异常入口未映射 & 故障地址、PDE/PTE、CR3 \\
中断后死机 & 栈、IDT、EOI 或清中断顺序错误 & SP、向量号、中断控制器状态 \\
DMA 后数据错 & cache 可见性、IOMMU、缓冲区所有权 & 描述符、completion、cache 维护记录 \\
\bottomrule
\end{longtable}

这张表把整本书的排查方法统一起来：不要先猜“CPU 坏了”，先问硬件行为停在哪个边界。电源和时钟是第一章的边界，触发器和状态是第二章的边界，ALU 和控制线是第三章的边界，PC/IR/寄存器是第四章的边界，主存/总线/MMIO/DMA 是第五章的边界，描述符、页表、固件和平台是第六章的边界。

\topic{案例：把教学板扩展成一台小系统}
最后给一个可以落地的整机扩展路线。第一阶段只跑 ROM 中的 LED 循环；第二阶段加入 SRAM 和栈，能调用函数；第三阶段加入定时器中断，能周期性切换状态；第四阶段加入简化 DMA，把一段内存复制到另一段；第五阶段加入保护或至少地址错误检测，让未映射访问进入错误入口。每一阶段都对应书中一个硬件层次。

\begin{longtable}{L{0.16\linewidth}L{0.42\linewidth}L{0.32\linewidth}}
\toprule
阶段 & 新增硬件/固件 & 预期行为 \\
\midrule
ROM LED & 复位向量、ROM、GPIO MMIO & 上电后固定节奏改写 LED 寄存器 \\
SRAM 栈 & SRAM 译码、SP 初始化、CALL/RET & 函数调用和局部变量不破坏全局数据 \\
定时器中断 & 定时器、向量表、中断清除 & 周期事件不丢失，主循环可继续运行 \\
简化 DMA & 描述符、源/目标地址、completion & DMA 完成后 CPU 读到正确内存内容 \\
错误入口 & 未映射检测、总线超时或简化异常 & 错误地址进入可观察处理程序 \\
\bottomrule
\end{longtable}

这条路线比“一步到位跑操作系统”更适合硬件书。读者每完成一层，都能指出新增了哪些地址、哪些控制线、哪些状态边界、哪些错误路径。等这些边界稳定后，再看 8086、80386 或现代 SoC，就不是背芯片史，而是在识别同一组问题如何被真实平台放大和工程化。

\topic{综合项目：写一份最小整机规格}
本章最后可以把全书内容收束成一份最小整机规格。规格不需要一开始很复杂，但必须让别人能按它搭出同一台机器，并能判断每个访问是成功、等待还是错误。下面给出一个建议模板，读者可以用分立逻辑、FPGA、微控制器外设模拟或软件仿真来实现。

\begin{longtable}{L{0.18\linewidth}L{0.42\linewidth}L{0.3\linewidth}}
\toprule
规格项目 & 必须写清楚 & 预期行为 \\
\midrule
地址空间 & ROM、RAM、MMIO、未映射区范围 & 任意地址最多一个目标响应 \\
复位行为 & PC 初值、SP 初值、控制线安全值 & 上电第一条取指可重复观察 \\
指令集 & 编码、寄存器字段、立即数扩展和标志规则 & 每条指令能写成控制 trace \\
内存访问 & 字节序、对齐、字节使能、等待和错误 & 加载/存储不破坏相邻字节 \\
设备寄存器 & 每个位的读写、副作用和清除规则 & GPIO/定时器/UART 状态可验证 \\
中断/异常 & 向量入口、保存状态、清除和返回规则 & 外部事件和错误能进入可诊断路径 \\
DMA/队列 & 描述符格式、所有权、completion 和 cache 规则 & CPU 与设备不同时拥有缓冲区 \\
\bottomrule
\end{longtable}

规格写完后，建议按下面的顺序实现，而不是同时接上所有模块。

\begin{longtable}{L{0.14\linewidth}L{0.42\linewidth}L{0.34\linewidth}}
\toprule
里程碑 & 要实现的最小能力 & 预期行为 \\
\midrule
M1 & 复位后从 ROM 取指，PC 顺序前进 & 逻辑分析仪能看到连续取指地址 \\
M2 & ADD/SUB/JNZ 正确执行计数循环 & 寄存器值和分支方向符合 trace \\
M3 & SRAM 加载/存储和栈可用 & 函数调用能返回，局部变量不乱 \\
M4 & GPIO MMIO 可读写 & LED/按键通过地址事务控制 \\
M5 & 定时器中断可进入和返回 & 中断不破坏主循环寄存器和栈 \\
M6 & 简化 DMA 或块复制控制器可搬数据 & completion 后内存内容正确 \\
M7 & 未映射访问进入错误入口 & 错误地址和原因可观察 \\
\bottomrule
\end{longtable}

每个里程碑都应保留一个最小测试程序。不要只保留最终大程序，因为最终程序一旦失败，无法判断是取指、译码、ALU、存储器、MMIO、中断还是 DMA 出错。小测试程序就是硬件书里的单元测试。

\begin{lstlisting}
M2 counter loop:
  MOV  R1, 0
  MOV  R2, 10
loop:
  ADD  R1, 1
  SUB  R2, 1
  JNZ  R2, loop
  MOV [GPIO_OUT], R1
\end{lstlisting}

\begin{longtable}{L{0.2\linewidth}L{0.38\linewidth}L{0.32\linewidth}}
\toprule
测试失败现象 & 优先怀疑 & 观察信号 \\
\midrule
PC 不变 & 复位、时钟、PC 写使能或取指等待 & \code{PC}、\code{reset}、\code{clock}、\code{ready} \\
循环次数错 & 立即数扩展、SUB、Zero 标志或 JNZ & \code{R2}、\code{Zero}、\code{pc_sel}、\code{branch_target} \\
GPIO 不变 & 地址译码、MMIO 写使能、字节使能 & \code{address}、\code{io_cs}、\code{mem_write}、\code{write_data} \\
函数返回错 & SP、CALL/RET 返回地址、栈 RAM & \code{SP}、栈内容、\code{PC}、\code{writeback} \\
中断后错乱 & 保存/恢复寄存器、EOI、清状态顺序 & \code{vector}、\code{SP}、saved PC、irq status \\
DMA 数据错 & 描述符、地址权限、cache 可见性、completion & \code{desc}、\code{buffer}、\code{owner}、\code{done} \\
\bottomrule
\end{longtable}

这个综合项目的目的，是让读者把书中的概念变成可执行机器：第一章的电平和输出驱动，第二章的状态边界，第三章的 ALU 和控制字，第四章的 CPU trace，第五章的地址事务和 DMA，第六章的启动、异常和平台约定。能写出并验证这份规格，才算真正把“硬件从零到整机”连起来。

\topic{综合项目的实验报告模板}
为了避免实验结果只停留在“好像能跑”，每个里程碑都应留下实验报告。报告不需要长篇叙事，但必须保存足够的观察记录，让另一个人能复现同样结论。建议每个实验至少记录以下字段：

\begin{longtable}{L{0.18\linewidth}L{0.42\linewidth}L{0.3\linewidth}}
\toprule
报告字段 & 要记录的内容 & 为什么重要 \\
\midrule
硬件版本 & 原理图版本、器件型号、时钟频率、电源电压 & 不同器件和频率会改变时序结论 \\
固件版本 & ROM 镜像、链接地址、入口地址、编译选项 & 同一硬件跑不同镜像可能表现不同 \\
地址地图 & ROM/RAM/MMIO/未映射区域范围 & 后续访问错误可直接对照 \\
测试程序 & 最小汇编或机器码列表 & 能定位具体指令和控制 trace \\
观察信号 & PC、IR、控制字、地址、数据、ready、写使能 & 证明结果来自正确硬件路径 \\
预期结果 & 预期寄存器、内存、设备和异常状态 & 避免把偶然现象当成通过 \\
失败记录 & 失败现象、假设、修改和复测结果 & 保留完整调试记录 \\
\bottomrule
\end{longtable}

例如 M3“SRAM 栈可用”的实验报告应包含：复位后 SP 的值，CALL 前后的 PC，栈上保存的返回地址，函数内部局部变量地址，RET 后 PC 是否回到调用点，R0 或指定返回值寄存器是否正确。若只写“函数调用成功”，后续出现异常时无法判断是栈、返回地址、寄存器保存还是取指路径出了问题。

\begin{lstlisting}
M3 trace example:
  reset PC = 0x00000000
  initial SP = 0x20008000
  CALL writes return address 0x00000124 to [SP-4]
  callee uses [SP-8] for local variable
  RET loads PC = 0x00000124
  R0 = expected return value
\end{lstlisting}

这类报告会把学习从“看懂章节”推进到“能维护机器”。硬件工程的核心不是一次点亮 LED，而是每次改动后仍能解释：哪个状态边界改变了，哪个事务完成了，哪个错误路径被覆盖了，哪些观察结果足以支持结论。

\topic{专题：从 8086 到 SoC 的“同一问题不同边界”}
8086、80386、Pentium、微控制器和现代 SoC 看起来差别很大，但平台工程反复面对同一组问题：第一条指令从哪里来，地址如何到达目标，慢设备如何等待，错误如何上报，外设如何通知 CPU，DMA 与 cache 如何保持可见，调试器如何观察状态。演进改变的是边界位置，而不是问题本身。

\begin{longtable}{L{0.18\linewidth}L{0.34\linewidth}L{0.38\linewidth}}
\toprule
平台问题 & 早期芯片形态 & 现代 SoC 形态 \\
\midrule
复位入口 & 外部 ROM 响应固定复位向量 & 片内 Boot ROM、熔丝、启动介质选择和安全检查 \\
地址译码 & 板级译码器产生片选 & 片上互连、地址防火墙、IOMMU 和外设桥 \\
等待状态 & READY/WAIT 拉长总线周期 & valid/ready、AXI response、QoS 和超时 \\
中断 & 外部中断控制器或简单向量 & GIC/APIC、优先级、亲和性、虚拟化中断 \\
DMA & 简单外设总线主控 & 多主设备、cache 一致性、SMMU/IOMMU 权限 \\
调试 & 逻辑分析仪看总线引脚 & JTAG/SWD、trace macrocell、性能计数器和固件日志 \\
\bottomrule
\end{longtable}

这张对照能帮助读者读芯片手册。看到 AXI、APB、GIC、IOMMU、Boot ROM、secure monitor 时，不应把它们当成孤立名词，而应问：它们分别把早期平台的哪一条裸露信号或板级功能搬进了芯片内部？搬进去之后，哪些行为变得更快、更安全，哪些调试信息变得更难直接观察？

\topic{专题：BIOS、Boot ROM 和 loader 的最小移交清单}
启动链路不是“跳到下一段代码”这么简单。每一级启动代码都必须向下一级交出可用的执行环境：当前特权级、栈、内存可用范围、异常入口、时钟、串口或日志设备、镜像位置和校验结果。缺少任何一项，下一阶段可能运行一段时间后才以难以定位的方式失败。

\begin{longtable}{L{0.2\linewidth}L{0.42\linewidth}L{0.28\linewidth}}
\toprule
移交项目 & 必须说明的内容 & 失败现象 \\
\midrule
入口地址 & 下一阶段第一条指令位置和执行状态 & 跳入错误模式或错误对齐地址 \\
栈 & 栈顶、增长方向、可用范围 & CALL/异常保存破坏数据或 ROM 区 \\
内存地图 & RAM、ROM、MMIO、保留区、安全区 & loader 覆盖设备寄存器或固件数据 \\
异常向量 & 向量基址、异常栈、默认处理器 & 早期错误无法诊断，直接死机 \\
时钟/电源 & CPU、总线、外设时钟是否打开 & 访问外设永久等待或超时 \\
日志通道 & UART、半主机、内存日志或调试口 & 失败时没有可观察的输出 \\
镜像验证 & 长度、哈希、签名、加载地址 & 执行损坏或不可信代码 \\
\bottomrule
\end{longtable}

安全启动只是把这张移交清单加上信任约束：Boot ROM 信任芯片内固化公钥或熔丝状态，验证下一阶段签名，再把控制权交给经过验证的镜像。若验证失败，平台不能继续执行任意代码，而应进入恢复、锁定或错误报告路径。这样读者能把“安全启动”理解为启动链路中的状态机和逐级验证过程，而不是抽象口号。

\topic{专题：平台 bring-up 日志怎样定位到具体层}
bring-up 的第一原则是每一步只引入一个新假设，并记录足够的观察结果。若一开始就运行完整系统，失败时无法知道是时钟、复位、内存、cache、异常、串口、设备树、驱动还是文件系统出错。更稳妥的办法是从复位向量开始逐层扩大可用机器。

\begin{longtable}{L{0.16\linewidth}L{0.42\linewidth}L{0.32\linewidth}}
\toprule
阶段 & 建议日志 & 证明的边界 \\
\midrule
复位 & PC、模式、第一条指令字 & ROM 映射和取指可用 \\
串口 & 时钟源、波特率、发送寄存器状态 & MMIO、外设时钟和基本输出可用 \\
SRAM & 写读 walking bit、地址线测试 & 片内 RAM 和数据访问可用 \\
DRAM & 初始化参数、训练结果、块读写校验 & 外部内存控制器和板级连线可用 \\
异常 & 故障地址、异常号、返回 PC & 向量表和异常栈可用 \\
cache/MMU & 页表摘要、属性、TLB flush 记录 & 地址转换和一致性边界可控 \\
中断 & 控制器配置、pending/active 状态 & 异步事件链路可用 \\
\bottomrule
\end{longtable}

每条日志都应能回答“这条记录来自哪里”。例如串口能打印字符，不只证明 C 函数被调用，还证明 CPU 能取指、写 MMIO、外设时钟打开、引脚复用正确、波特率近似正确。若后续打开 cache 后打印乱码，就可以把变化集中到 cache 属性、写缓冲、MMIO ordering 或时钟切换，而不是重新怀疑所有层。

\topic{专题：把经典芯片放进同一张学习路线}
本章不要求记住每颗芯片的全部引脚，而要求形成可迁移的读法。读 6502 时看复位向量、零页和简单总线；读 8086 时看段:偏移、复用总线和字节存储体；读 80386 时看保护模式、分页和异常；读 Pentium 时看流水线、cache 和总线协议；读 SoC 时看片上互连、启动 ROM、DMA、一致性和安全边界。

\begin{longtable}{L{0.18\linewidth}L{0.42\linewidth}L{0.3\linewidth}}
\toprule
阅读对象 & 优先抓住的问题 & 可迁移能力 \\
\midrule
6502/Z80/8051 & 复位、寄存器、地址空间、简单外设 & 从最少状态读懂机器执行 \\
8086 & 分段、外部总线、READY、板级译码 & 把 CPU 周期映射到平台信号 \\
80386 & 特权级、描述符、分页、异常 & 理解操作系统为什么需要硬件保护 \\
80486/Pentium & cache、流水线、总线事务 & 解释性能和一致性不再只是 ISA 问题 \\
现代 SoC & Boot ROM、片上互连、DMA、安全、调试 & 把整机看成多主设备协作系统 \\
\bottomrule
\end{longtable}

这样安排后，经典芯片不是怀旧材料，而是复杂系统的分层样本。每前进一步都在回答同一件事：更多性能、更多外设、更多安全约束进入系统后，哪些状态必须被显式保存，哪些事务必须被确认完成，哪些错误必须变成可诊断信息。

\topic{专题：主板不是插槽集合，而是平台约定}
主板把 CPU、内存、固件、电源、时钟、PCIe 设备、USB/SATA/NVMe、显卡、网卡和低速管理控制器连成一个可启动平台。它的核心职责不是“提供很多接口”，而是保证复位后每个设备处于可枚举、可供电、可配置、可报告错误的状态。CPU 能跑程序只是平台的一部分；内存训练、PCIe 枚举、电源时序、时钟分配、固件配置和热管理同样决定整机是否可靠。

\begin{pdfdiagram}
  \node[hw state, minimum width=1.55cm] (cpu) at (0,0.8) {CPU};
  \node[hw memory, minimum width=1.75cm] (dram) at (2.6,1.65) {DDR 内存};
  \node[hw compute, minimum width=1.9cm] (pch) at (2.6,-0.35) {芯片组/PCH};
  \node[hw control, minimum width=1.7cm] (fw) at (0,-1.35) {UEFI/Flash};
  \node[hw control, minimum width=1.65cm] (pcie) at (5.2,0.8) {PCIe Root};
  \node[hw state, minimum width=1.55cm] (gpu) at (7.5,1.55) {显卡};
  \node[hw control, minimum width=1.55cm] (nvme) at (7.5,0.15) {NVMe};
  \node[hw control, minimum width=1.55cm] (usb) at (5.2,-1.35) {USB/SATA/网卡};
  \draw[hw bus] (cpu) -- node[above,hw label,fill=white]{内存控制器} (dram);
  \draw[hw bus] (cpu) -- node[left,hw label,fill=white]{DMI/片上互连} (pch);
  \draw[hw bus] (fw) -- node[left,hw label,fill=white]{启动固件} (pch);
  \draw[hw bus] (cpu) -- (pcie);
  \draw[hw bus] (pcie) -- (gpu);
  \draw[hw bus] (pcie) -- (nvme);
  \draw[hw bus] (pch) -- (usb);
\end{pdfdiagram}
\diagramnote{PC 主板的关键不是“很多线”，而是启动、枚举、供电、时钟和错误报告共同构成平台约定。}

\begin{longtable}{L{0.18\linewidth}L{0.36\linewidth}L{0.36\linewidth}}
\toprule
主板对象 & 负责的问题 & 典型失败表现 \\
\midrule
供电 VRM & 把电源转换成 CPU/GPU/内存所需电压，并按时序上电 & 复位不释放、负载变化时掉压、机器重启 \\
时钟源 & 给 CPU、内存、PCIe、USB 等提供参考时钟 & 链路训练失败、设备枚举不稳定 \\
固件 Flash & 保存 UEFI/BIOS、微码、平台配置 & 无法取第一阶段代码或配置错误 \\
DDR 通道 & 连接内存颗粒或 DIMM，完成训练和 ECC & 内存训练失败、随机数据错误 \\
PCIe Root & 枚举显卡、NVMe、网卡和扩展设备 & 设备消失、降速、AER 错误 \\
低速管理 & Super I/O、EC、传感器、风扇、RTC & 键盘、风扇、温度或睡眠唤醒异常 \\
\bottomrule
\end{longtable}

\topic{专题：PCIe 是分组互连，不是并行总线升级版}
早期外部总线常暴露地址线、数据线、读写控制和 READY；PCIe 则把事务封装成包，在点对点串行链路上传输。CPU 或芯片组通过 root complex 发起配置、内存读写和消息事务；设备通过 BAR 暴露 MMIO 窗口，通过 MSI/MSI-X 发送中断，通过 DMA 访问主存。PCIe 的关键边界包括链路训练、配置空间、BAR 映射、事务层包、完成包和错误报告。

\begin{longtable}{L{0.18\linewidth}L{0.38\linewidth}L{0.34\linewidth}}
\toprule
PCIe 概念 & 含义 & 本书对应概念 \\
\midrule
配置空间 & 枚举设备、读取 ID、配置 BAR 和能力 & 地址地图和平台发现 \\
BAR & 把设备寄存器映射到 CPU 地址空间 & MMIO 副作用和访问属性 \\
Memory Read/Write TLP & 设备或 root 发起的内存事务包 & 总线事务和 completion \\
MSI/MSI-X & 设备通过写内存样式消息触发中断 & 中断不一定是独立引脚 \\
AER & 高级错误报告 & 错误路径可诊断 \\
\bottomrule
\end{longtable}

把 PCIe 讲清楚后，显卡、NVMe、网卡都可以归入同一模型：先由平台枚举设备并分配地址窗口；驱动把命令写入队列或 MMIO；设备通过 DMA 读写主存；完成后写 completion 或发 MSI；错误通过状态寄存器、completion 状态或 AER 上报。这样读者不会把“插在主板上的设备”理解成 CPU 的附属品，而会把它看成能主动访问内存的总线主设备。

\topic{专题：显卡的核心不是显示器接口，而是并行处理器加显存系统}
显卡既是显示设备，也是大规模并行计算设备。GPU 内部有许多执行单元、调度器、寄存器文件、L1/L2 cache、纹理/采样单元、光栅化和显示引擎；外部或封装上有高带宽显存，例如 GDDR 或 HBM。CPU 不会逐像素把画面写到显示器，而是准备命令缓冲区、资源描述、纹理和顶点数据，GPU 读取这些对象并在显存中生成帧缓冲，显示引擎再按扫描时序把帧送到显示接口。

\begin{pdfdiagram}
  \node[hw state, minimum width=1.6cm] (cpu) at (0,0) {CPU};
  \node[hw memory, minimum width=2.0cm] (sysmem) at (2.4,0.8) {系统内存};
  \node[hw control, minimum width=2.0cm] (cmd) at (2.4,-0.8) {命令缓冲};
  \node[hw control, minimum width=1.85cm] (pcie) at (4.9,0) {PCIe};
  \node[hw compute, minimum width=2.0cm] (gpu) at (7.35,0.8) {GPU 执行};
  \node[hw memory, minimum width=1.75cm] (vram) at (9.7,0.8) {VRAM};
  \node[hw state, minimum width=1.75cm] (scan) at (7.35,-0.8) {显示引擎};
  \node[hw state, minimum width=1.55cm] (disp) at (9.7,-0.8) {显示器};
  \draw[hw bus] (cpu) -- (sysmem);
  \draw[hw bus] (cpu) -- (cmd);
  \draw[hw bus] (cmd) -- (pcie);
  \draw[hw bus] (pcie) -- (gpu);
  \draw[hw bus] (gpu) -- (vram);
  \draw[hw return] (vram) -- (scan);
  \draw[hw bus] (scan) -- (disp);
\end{pdfdiagram}
\diagramnote{GPU 数据路径：CPU 提交命令，GPU 经 PCIe/内存读取资源，在显存中生成帧，显示引擎独立扫描输出。}

\begin{longtable}{L{0.18\linewidth}L{0.38\linewidth}L{0.34\linewidth}}
\toprule
显卡对象 & 硬件含义 & 常见误解 \\
\midrule
命令缓冲 & CPU 写入、GPU 异步消费的工作队列 & 不是 CPU 直接调用 GPU 函数 \\
VRAM & GPU 近处的大带宽存储系统 & 不是普通 DRAM 的别名，带宽和访问模式不同 \\
着色/计算单元 & 大量相似线程并行执行 & 不是少数超快 CPU 核心 \\
纹理/cache 单元 & 优化二维或三维局部访问和采样 & 不等于通用 cache 的简单复制 \\
显示引擎 & 按固定时序读取帧缓冲并输出信号 & 渲染完成和显示扫描是两个边界 \\
\bottomrule
\end{longtable}

GPU 的性能问题也不同于 CPU。CPU 常追求低时延和强单线程控制流；GPU 常追求吞吐，用大量并行线程隐藏显存等待。分支发散、访存不合并、显存带宽不足、CPU/GPU 同步过多，都会让 GPU 性能掉下来。理解显卡时应问：命令何时提交，资源何时对 GPU 可见，GPU 何时完成，帧缓冲何时可显示，CPU 是否过早读回结果。

\topic{专题：整机案例：从打开程序到读取文件和显示图像}
一个真实用户动作会穿过多种硬件。用户打开程序，操作系统可能从 SSD 读取可执行文件和页面；CPU 触发块设备队列；NVMe 控制器 DMA 数据到主存；页表和 cache 让 CPU 取指执行；程序提交图形命令；GPU 读取资源、渲染到 VRAM；显示引擎扫描帧缓冲。这里没有哪个部件可以单独解释全程，必须把 CPU、内存、SSD、PCIe、GPU、显示和中断连成 trace。

\begin{longtable}{L{0.16\linewidth}L{0.44\linewidth}L{0.3\linewidth}}
\toprule
阶段 & 硬件动作 & 完成标志 \\
\midrule
缺页/读文件 & OS 提交 NVMe 读请求，设备 DMA 到主存 & 块 I/O completion 和页状态有效 \\
执行代码 & CPU 取指、TLB/cache 查找、分支和写回 & 指令提交，异常路径为空 \\
准备图形资源 & CPU 写命令缓冲和资源描述符 & 命令队列对 GPU 可见 \\
GPU 渲染 & GPU 读取资源，执行 shader，写 VRAM & fence 或 completion 表示渲染完成 \\
显示扫描 & 显示引擎按刷新时序读取帧缓冲 & vblank、翻页或显示控制器状态 \\
\bottomrule
\end{longtable}

这个案例可以作为读完整本书后的纸面项目。要求不是实现操作系统和驱动，而是把每一步的机器行为写清楚：谁拥有缓冲区，谁发起 DMA，哪个地址空间有效，cache/TLB 在哪里参与，哪个 completion 才说明完成，哪个错误会阻止下一阶段。

\topic{专题：纸面项目应替代不可执行的大实验}
如果没有硬件平台，仍然可以做高质量项目。项目形式应从“搭出来能跑”改成“规格写得可检查”。例如给出一台简化 PC：一个 CPU、两级 cache、DDR 控制器、PCIe root、NVMe、GPU、USB 控制器和固件 ROM。读者要写地址地图、启动顺序、设备枚举、一次文件读取、一次 GPU 命令提交、一次中断处理和一个错误路径。这样的纸面项目比空泛地说“了解主板”更有效。

\begin{longtable}{L{0.2\linewidth}L{0.46\linewidth}L{0.24\linewidth}}
\toprule
纸面项目 & 必须交付 & 合格标准 \\
\midrule
简化 PC 地址地图 & DRAM、固件、MMIO、PCIe BAR、未映射区 & 任意地址有明确目标或错误 \\
NVMe 读路径 & submission、doorbell、DMA buffer、completion、中断 & CPU 不把门铃返回当成 I/O 完成 \\
GPU 命令路径 & 命令缓冲、资源可见性、VRAM、fence、显示扫描 & 渲染完成和显示完成分开 \\
主板启动路径 & 电源、时钟、复位、固件、内存训练、PCIe 枚举 & 每一步有失败后可诊断状态 \\
错误注入 & TLB miss、cache miss、PCIe 错误、SSD 超时、GPU hang & 错误归属到具体层，不混成“程序卡住” \\
\bottomrule
\end{longtable}



\section*{章末练习}

\begin{longtable}{L{0.2\linewidth}L{0.5\linewidth}L{0.2\linewidth}}
\toprule
任务 & 要求 & 学习目标 \\
\midrule
8086 地址形成 & 用段寄存器和偏移计算一个物理地址 & 能说明 20-bit 地址来自硬件加法路径 \\
复用总线 & 写出 8086 读周期中地址、ALE、数据方向和 READY 的阶段 & 能指出何时锁存地址、何时采样数据 \\
字节存储体 & 解释 A0 与 \(\overline{\mathrm{BHE}}\) 如何选择低/高字节 & 能处理奇地址字访问 \\
80386 分页 & 把一个 32-bit 线性地址拆成目录索引、页表索引和页内偏移 & 能说明异常发生在真正访存前 \\
保护模式入口 & 列出 GDTR/IDTR、CR0、CR3 和跳转刷新各自作用 & 能说明状态切换顺序 \\
复位向量 & 比较 6502、Z80/8051、8086、80386 后续 x86 的复位入口 & 能说明第一条指令必须由非易失存储响应 \\
平台分层 & 说明 CPU 本地总线、内存/cache 控制、高速扩展、低速 I/O 和固件层职责 & 能说明等待来自整个平台 \\
SoC 演进 & 说明把内存控制器、I/O 和安全启动搬进片内的收益与代价 & 能指出片上互连和调试复杂度 \\
主板结构 & 画出 CPU、DDR、PCH、UEFI Flash、PCIe、NVMe、GPU 的连接关系 & 能说明供电、时钟、复位和枚举不是 CPU 内部问题 \\
PCIe 设备 & 说明配置空间、BAR、TLP、MSI 和 AER 各自解决什么问题 & 能把显卡/NVMe/网卡统一成总线主设备 \\
显卡路径 & 写出 CPU 提交命令、GPU 读资源、VRAM 写帧、显示扫描的 trace & 能区分渲染完成、DMA 完成和显示完成 \\
6502 寻址与拍数 & 用零页寻址和 (zp),Y 间接变址写数组求和循环并逐拍计数 & 理解拍数就是总线事务序列 \\
Z80 交换与刷新 & 说明 EXX 怎样加速中断响应、R 寄存器怎样把刷新藏进取指 & 寄存器组交换=不访存的上下文切换 \\
486/Pentium 集成 & 列出 80486 和 Pentium 各自搬进片内的部件及对应收益 & 集成把等待从板级移到片内 \\
\bottomrule
\end{longtable}

\section*{参考解答与要点提示}

\begin{longtable}{L{0.18\linewidth}L{0.5\linewidth}L{0.22\linewidth}}
\toprule
任务 & 参考解答 & 必须掌握的要点 \\
\midrule
8086 地址形成 & 物理地址由 \code{segment*16+offset} 得到，内部编程模型保持 16-bit，外部地址能力扩展到 20-bit。 & 分段是地址形成硬件。 \\
复用总线 & T1 输出地址并用 ALE 锁存；T2 改变总线方向并给出读控制；T3/Tw 等待目标驱动数据且 READY 可延长；T4 采样并结束。 & 复用引脚必须分阶段解释。 \\
字节存储体 & 偶地址低字节由 A0 选择，奇地址高字节由 \(\overline{\mathrm{BHE}}\) 参与选择；奇地址字访问可能拆成两个周期。 & 位宽不等于任意地址一次完成。 \\
80386 分页 & 4 KiB 页下，高 10 bit 选 PDE，中 10 bit 选 PTE，低 12 bit 是页内偏移；PDE/PTE 不存在或权限不允许时触发异常。 & 地址转换也是硬件路径。 \\
保护模式入口 & 进入保护模式前要准备描述符表、装载 GDTR/IDTR、设置 CR0、刷新取指路径；分页还需准备页目录/页表、装载 CR3、打开分页位。 & 状态切换顺序错误会让取指和异常入口失效。 \\
复位向量 & 6502 从高端向量取入口，Z80/8051 常从 0 地址开始，8086 从 \code{0xFFFF0} 附近，后续 x86 从更高物理地址附近，SoC 可从 Boot ROM 进入。 & 启动入口是硬件规格。 \\
平台分层 & CPU 本地总线处理核心事务，内存/cache 控制处理 DRAM 和回填，高速扩展处理图形/磁盘/网络，低速 I/O 处理串口/键盘/RTC，固件负责初始化和启动选择。 & 主频之外还有平台等待。 \\
SoC 演进 & 片上集成缩短常见路径、减少引脚和功耗，但引入片上互连、一致性、时钟/电源域、安全启动和调试复杂度。 & 低时延来自边界移动，不是所有路径变快。 \\
主板结构 & CPU 经前端总线连北桥，北桥负责内存和图形，南桥挂低速 I/O（串口、键盘、RTC），固件 Flash 上电先执行并完成初始化。 & 平台分层是把不同速度的路径分开组织。 \\
PCIe 设备 & 固件读配置空间枚举设备并分配 BAR 地址窗口，驱动经 MMIO 配置队列，设备用 DMA 搬数据，MSI/MSI-X 投递完成中断，AER 报告错误。 & 枚举-配置-DMA-完成是统一模型。 \\
显卡路径 & CPU 写命令环并按 doorbell，GPU 读资源和命令渲染，结果写 VRAM，显示控制器按扫描时序读出；渲染完成、DMA 完成、显示完成是三件不同的事。 & 三种完成语义不能混。 \\
6502 寻址与拍数 & LDA 零页 3 拍，(zp),Y 间接变址 5 拍（跨页 6 拍）；求和循环每轮用 (zp),Y 读一个元素，接 INY、BNE，拍数可手工核对。 & 简单数据通路让逐拍 trace 可行。 \\
Z80 交换与刷新 & EXX 一条指令交换全部通用寄存器，省去压栈；M1 取指后半段把 R 输出到地址总线完成 DRAM 刷新。 & 交换和借用都是减少总线事务。 \\
486/Pentium 集成 & 486 搬入 cache、FPU、流水线；Pentium 搬入双发射、BTB 预测、分离 cache 和 64-bit 数据总线。 & 每一代都在移动等待的位置。 \\
\bottomrule
\end{longtable}
